% Worked solutions for Listas/lista-r5b-trade-credit-working-capital.md (15 exercises).
% Compile: pdflatex twice; check: pdffonts (Type 1 only, uni = yes).
% Every number below is asserted in lista-r5b-trade-credit-working-capital-solutions_check.py.
\documentclass[11pt,a4paper]{article}
\usepackage[utf8]{inputenc}
\usepackage[T1]{fontenc}
\usepackage{lmodern}   % Type 1 vector fonts (replaces PK bitmaps)
\usepackage{cmap}      % ToUnicode CMap -> copyable text and accents
\usepackage[brazilian,english]{babel}
\usepackage{amsmath,amssymb,amsfonts}
\usepackage[a4paper,margin=2.2cm]{geometry}
\usepackage{booktabs}
\usepackage{enumitem}
\usepackage{fancyhdr}
\usepackage{xcolor}
\usepackage{microtype}

\pagestyle{fancy}\fancyhf{}
\lhead{\footnotesize Finan\c{c}as II --- Corporativas (MPE Insper)}
\rhead{\footnotesize List R5b --- Worked solutions}
\cfoot{\footnotesize\thepage}
\renewcommand{\headrulewidth}{0.3pt}
\setlength{\headheight}{14pt}
\setlist{topsep=2pt,itemsep=1pt}

\newcommand{\exsec}[1]{\section*{\color{blue!55!black}#1}\addcontentsline{toc}{section}{#1}}
\newcommand{\notation}{\par\noindent\textbf{Notation.}}
\newcommand{\step}[2]{\par\medskip\noindent\textbf{Step #1 --- #2.}~}
\newcommand{\intuition}{\par\medskip\noindent\textbf{Intuition.}~}
\newcommand{\bookans}{\par\medskip\noindent\textbf{Book answer.}~}
\newcommand{\refbib}{\par\smallskip\noindent\textbf{Ref:}~}
\newcommand{\assume}{\par\smallskip\noindent\textbf{Assumption.}~}
\newcommand{\pt}{\foreignlanguage{brazilian}}
\newcommand{\asked}{\par\medskip\noindent\textbf{What is asked.}~}
\newcommand{\approach}{\par\smallskip\noindent\textbf{Approach.}~}
\newcommand{\answer}[1]{\par\smallskip\noindent\fbox{\textbf{Answer #1.}}~}
\newcommand{\pitfalls}{\par\medskip\noindent\textbf{Common mistakes.}}

\title{List R5b --- Trade credit and working capital management\\[2pt]
\large Worked solutions}
\author{Finan\c{c}as II --- Corporativas (MPE Insper)}
\date{}

%% STATEMENT-PREAMBLE BEGIN (tools/statements_into_solutions.py)
\usepackage{longtable,array,calc}
\providecommand{\tightlist}{\setlength{\itemsep}{0pt}\setlength{\parskip}{0pt}}
\providecommand{\href}[2]{#2}
\providecommand{\real}[1]{#1}
\makeatletter\@ifundefined{c@none}{\newcounter{none}}{}\makeatother  % pandoc's \LTcaptype{none}
\newcommand{\stmtbegin}{\par\medskip\noindent\textbf{Statement.}\par\nopagebreak\begingroup\small\setlength{\tabcolsep}{3pt}\sloppy}
\newcommand{\stmtend}{\par\endgroup\nopagebreak\noindent\rule{\linewidth}{0.4pt}\par\medskip}
\DeclareUnicodeCharacter{2192}{\ensuremath{\rightarrow}}
\DeclareUnicodeCharacter{2212}{\ensuremath{-}}
\DeclareUnicodeCharacter{2022}{\textbullet{}}
\DeclareUnicodeCharacter{2248}{\ensuremath{\approx}}
\DeclareUnicodeCharacter{2032}{\ensuremath{^{\prime}}}
\DeclareUnicodeCharacter{2081}{\ensuremath{_{1}}}
\DeclareUnicodeCharacter{03B8}{\ensuremath{\theta}}
\DeclareUnicodeCharacter{03BB}{\ensuremath{\lambda}}
\DeclareUnicodeCharacter{03C8}{\ensuremath{\psi}}
\DeclareUnicodeCharacter{03B1}{\ensuremath{\alpha}}
\DeclareUnicodeCharacter{25A0}{\rule{0.7em}{0.7em}}
%% STATEMENT-PREAMBLE END
\begin{document}
\maketitle
\noindent Statements are in \texttt{Listas/lista-r5b-trade-credit-working-capital.md}; each
exercise is cited by number and printed page and its full text is reproduced right before its solution. Editions are the local
copies: Berk = Berk, DeMarzo \& Harford, \emph{Fundamentals of Corporate Finance}, 2nd ed.
(2011); Ross = Ross, Westerfield, Jordan \& Lamb, \emph{Fundamentos de Administra\c{c}\~ao
Financeira}, 9th ed. (2013); Higgins = Koski, Higgins \& Mitton, \emph{Analysis for Financial
Management}, 13th ed. (2023); Grinblatt = Grinblatt \& Titman, \emph{Financial Markets and
Corporate Strategy}, 2nd ed. (2002). Ross prints thousands with a period (\pt{\$89.804}); below
they are written $89{,}804$. Money is in the units of each statement.

\medskip\noindent\textbf{Conventions used throughout.} A year has 365 days (all four books use
365). Days ratios use the year-flow divided by 365: inventory and payables days on cost of goods
sold (COGS), receivables days on sales. The effective annual rate (EAR) of a rate $r$ earned over
$n$ days is $(1+r)^{365/n}-1$, i.e.\ compounding over $365/n$ periods (Berk Eq.~5.1; Ross,
p.~709).

\tableofcontents

% =====================================================================
\exsec{1. Westerly Industries: cash conversion cycle --- Berk 19.3, p.~586}
%% BEGIN-STATEMENT 1
\stmtbegin
\textbf{Source:} Berk, Ch. 19, Exercise 3, p.~586 (PDF 617) · Answer in book: no

Westerly Industries has the following financial information. What is its cash conversion cycle?

{\def\LTcaptype{none} % do not increment counter
\begin{longtable}[]{@{}ll@{}}
\toprule\noalign{}
Item & Amount \\
\midrule\noalign{}
\endhead
\bottomrule\noalign{}
\endlastfoot
Sales & 100,000 \\
Cost of Goods Sold & 80,000 \\
Accounts Receivable & 30,000 \\
Inventory & 15,000 \\
Accounts Payable & 40,000 \\
\end{longtable}
}
\stmtend
%% END-STATEMENT

\asked Compute the cash conversion cycle (CCC): the number of days between paying suppliers for
inventory and collecting cash from customers, from one year of sales and COGS and the year-end
balances of receivables, inventory and payables.

\approach The CCC is the sum of inventory days and receivables days (the operating cycle) minus
payables days (Berk, \S19.1). Each ``days'' figure converts a balance into the number of days of
the matching flow it represents: inventory and payables are carried at cost, so they are divided
by daily COGS; receivables are carried at selling price, so they are divided by daily sales.

\notation
\begin{itemize}
  \item $S$, $C$: annual sales and COGS (\$); $AR$, $INV$, $AP$: receivables, inventory, payables (\$).
  \item $D_{INV}$, $D_{AR}$, $D_{AP}$: inventory, receivables and payables days (days, $\ge 0$).
  \item $OC = D_{INV}+D_{AR}$: operating cycle; $CCC = OC - D_{AP}$ (days, any sign).
\end{itemize}
\assume The balances are representative of the year (no averages are given).

\step{1}{Daily flows}
Divide each annual flow by 365 to get the denominators of the days ratios:
\begin{align}
  C/365 &= (80{,}000)/365 = 219.18, \notag\\
  S/365 &= (100{,}000)/365 = 273.97. \notag
\end{align}

\step{2}{The three days ratios}
\begin{align}
  D_{INV} &= \frac{INV}{C/365} = \frac{(15{,}000)}{(80{,}000/365)} = 68.44 \ \text{days}, \label{r5be1-inv}\\
  D_{AR} &= \frac{AR}{S/365} = \frac{(30{,}000)}{(100{,}000/365)} = 109.50 \ \text{days}, \label{r5be1-ar}\\
  D_{AP} &= \frac{AP}{C/365} = \frac{(40{,}000)}{(80{,}000/365)} = 182.50 \ \text{days}. \label{r5be1-ap}
\end{align}

\step{3}{Cash conversion cycle}
Add inventory and receivables days, then subtract payables days:
\begin{align}
  CCC &= D_{INV} + D_{AR} - D_{AP} \label{r5be1-ccc}\\
  &= (68.44) + (109.50) - (182.50) \notag\\
  &= (177.94) - (182.50) \notag\\
  &= -4.56 \ \text{days}. \notag
\end{align}

\answer{\unskip} $CCC = -4.56$ days (operating cycle $177.94$ days, payables $182.50$ days).

\intuition A negative CCC means suppliers wait longer to be paid (half a year) than it takes
Westerly to turn inventory into cash; suppliers finance the whole operating cycle, and a little
more. The firm needs no cash of its own to carry working capital --- the opposite of the usual
case. Whether this is good depends on why: if Westerly is forgoing early-payment discounts or
stretching beyond the due date, the ``free'' financing may be very expensive (Exercises 7 and 13).

\pitfalls
\begin{itemize}
  \item Dividing payables or inventory by sales instead of COGS: it mixes cost-based balances with a
        price-based flow and understates both ratios.
  \item Reporting the operating cycle ($177.94$ days) as the answer, i.e.\ forgetting to net out payables.
  \item Treating the negative sign as an arithmetic error; it is a legitimate outcome.
\end{itemize}

\refbib Berk Ch.~19, \S19.1 ``Overview of Working Capital'', The Cash Cycle, p.~567. No printed answer.

% =====================================================================
\exsec{2. Mediadora: operating and cash cycles --- Ross 18.6, p.~654}
%% BEGIN-STATEMENT 2
\stmtbegin
\textbf{Source:} Ross, Cap. 18, Exercise 18.6, p.~654 (PDF 678) · Answer in book: no

\textbf{Cálculo dos ciclos} {[}OA1{]} Considere as seguintes informações da demonstração contábil da
Mediadora S/A:

{\def\LTcaptype{none} % do not increment counter
\begin{longtable}[]{@{}lll@{}}
\toprule\noalign{}
Item & Inicial & Final \\
\midrule\noalign{}
\endhead
\bottomrule\noalign{}
\endlastfoot
Estoque & \$9.780 & \$11.380 \\
Contas a receber & 4.108 & 4.938 \\
Contas a pagar & 7.636 & 7.927 \\
Vendas a prazo & \$89.804 & \\
Custo das mercadorias vendidas & 56.384 & \\
\end{longtable}
}

Calcule os ciclos operacional e financeiro. Como você interpreta sua resposta?
\stmtend
%% END-STATEMENT

\asked Compute the operating cycle (days from buying inventory to collecting the sale) and the cash
cycle (days from paying suppliers to collecting the sale) from beginning and ending balances,
credit sales and COGS, and interpret them.

\approach Ross (\S18.2) computes each period from a turnover ratio built on the \emph{average}
of beginning and ending balances: inventory and payables turnover on COGS, receivables turnover
on credit sales; days $=365/\text{turnover}$. Operating cycle $=$ inventory period $+$
receivables period; cash cycle $=$ operating cycle $-$ payables period.

\notation
\begin{itemize}
  \item $\overline{INV}$, $\overline{AR}$, $\overline{AP}$: average inventory, receivables, payables (\$).
  \item $T_{INV}$, $T_{AR}$, $T_{AP}$: turnovers (times per year); $D_{INV}$, $D_{AR}$, $D_{AP}$: periods (days).
  \item $S_c = 89{,}804$: credit sales; $C = 56{,}384$: COGS (\$).
\end{itemize}

\step{1}{Average balances}
\begin{align}
  \overline{INV} &= (9{,}780 + 11{,}380)/2 = 10{,}580, \notag\\
  \overline{AR} &= (4{,}108 + 4{,}938)/2 = 4{,}523, \notag\\
  \overline{AP} &= (7{,}636 + 7{,}927)/2 = 7{,}781.5. \notag
\end{align}

\step{2}{Turnovers}
\begin{align}
  T_{INV} &= C/\overline{INV} = (56{,}384)/(10{,}580) = 5.3293, \notag\\
  T_{AR} &= S_c/\overline{AR} = (89{,}804)/(4{,}523) = 19.8550, \notag\\
  T_{AP} &= C/\overline{AP} = (56{,}384)/(7{,}781.5) = 7.2459. \notag
\end{align}

\step{3}{Periods in days}
\begin{align}
  D_{INV} &= 365/T_{INV} = 365/(5.3293) = 68.49 \ \text{days}, \label{r5be2-inv}\\
  D_{AR} &= 365/T_{AR} = 365/(19.8550) = 18.38 \ \text{days}, \label{r5be2-ar}\\
  D_{AP} &= 365/T_{AP} = 365/(7.2459) = 50.37 \ \text{days}. \label{r5be2-ap}
\end{align}

\step{4}{The two cycles}
\begin{align}
  OC &= D_{INV} + D_{AR} = (68.49) + (18.38) = 86.87 \ \text{days}, \label{r5be2-oc}\\
  CC &= OC - D_{AP} = (86.87) - (50.37) = 36.50 \ \text{days}. \label{r5be2-cc}
\end{align}

\answer{\unskip} Operating cycle $= 86.87$ days; cash cycle $= 36.50$ days.

\noindent\textbf{Interpretation.} Goods sit in inventory for about 68 days and customers take
about 18 days to pay, so it takes 87 days to turn a purchase into cash. Suppliers finance 50 of
those days. For the remaining $36.5$ days the firm has paid out cash that has not yet come back,
and that gap must be financed (cash, a bank line or equity). Most of the cycle is inventory, so
inventory management is where days can be cut; receivables are already quick.

\intuition The cash cycle measures how long the firm's own money is tied up per unit of
activity. Multiplied by daily COGS ($56{,}384/365 \approx 154$ per day) it gives an order of
magnitude for the working capital that must be funded; any day saved frees that much cash.

\pitfalls
\begin{itemize}
  \item Using ending instead of average balances; Ross uses averages when both are given.
  \item Building receivables turnover on COGS or payables turnover on sales.
  \item Interpreting the cash cycle as a profitability measure; it is a financing-need measure.
\end{itemize}

\refbib Ross Cap.~18, \S18.2 ``O ciclo operacional e o ciclo financeiro'', pp.~627--631 (worked
example in the chapter review, 18.1, pp.~651--652). No printed answer.

% =====================================================================
\exsec{3. Cash discounts on a 400-unit order --- Ross 20.1, p.~733}
%% BEGIN-STATEMENT 3
\stmtbegin
\textbf{Source:} Ross, Cap. 20, Exercise 20.1, p.~733 (PDF 757) · Answer in book: no

\textbf{Descontos financeiros} {[}OA1{]} Você faz um pedido de 400 unidades para estoque a um preço
unitário de \$125. O fornecedor oferece prazo de pagamento de 30 dias e desconto de 1\% para
pagamento em até 10 dias.

\begin{enumerate}
\def\labelenumi{\alph{enumi}.}
\item
  Quanto tempo você tem para pagar até que a conta vença? Se você aceitar o período completo, quanto
  deverá pagar?
\item
  Qual é o desconto oferecido? Em quanto tempo você deve pagar para obter o desconto? Se você
  aceitar o desconto, quanto pagará?
\item
  Se você não aceitar o desconto, quanto de juros pagará implicitamente? Quantos dias de crédito
  você está recebendo?
\end{enumerate}
\stmtend
%% END-STATEMENT

\asked (a) How long until the bill is due and how much is paid on the due date. (b) The size of
the discount, the deadline to get it, and the discounted payment. (c) How much interest is
implicitly paid by skipping the discount, and for how many days of credit.

\approach The terms ``1/10, net 30'' mean: pay the full invoice by day 30, or pay 1\% less by day
10. Skipping the discount is a loan from the supplier of the \emph{discounted} amount for the
days between the discount deadline and the due date; the forgone discount is the interest (Ross,
\S20.2, ``Custo de cr\'edito'').

\notation
\begin{itemize}
  \item $F$: invoice amount (\$); $d = 0.01$: discount rate; $t_d = 10$, $t_n = 30$: discount and net days.
  \item $r$: interest rate for the $t_n - t_d$ days of credit; EAR: its effective annual equivalent.
\end{itemize}

\step{1}{Invoice amount}
$F = (400)(125) = 50{,}000$.

\step{2}{(a) Full credit period}
The bill is due on day $t_n = 30$; paying then costs the full $F = \$50{,}000$.

\answer{(a)} 30 days; pay \$50{,}000.

\step{3}{(b) Discount}
\begin{align}
  dF &= (0.01)(50{,}000) = 500, \notag\\
  F - dF &= (50{,}000) - (500) = 49{,}500. \notag
\end{align}

\answer{(b)} A 1\% discount (\$500) for paying within 10 days; pay \$49{,}500.

\step{4}{(c) Implicit interest and days of credit}
Not taking the discount means keeping \$49{,}500 for $30-10 = 20$ extra days and paying \$500
more. The rate for those 20 days and its annual equivalent are
\begin{align}
  r &= \frac{dF}{F - dF} = \frac{(500)}{(49{,}500)} = 1.0101\%, \label{r5be3-r}\\
  \text{EAR} &= (1+r)^{365/20} - 1 = (1.010101)^{18.25} - 1 = 20.13\%. \label{r5be3-ear}
\end{align}

\answer{(c)} Implicit interest of \$500 for 20 days of credit ($1.0101\%$ per 20 days, an EAR of $20.13\%$).

\intuition The discount looks tiny, but it is earned for paying only 20 days early; repeated 18
times a year it compounds to about 20\%. A buyer that can borrow below 20\% should borrow and take
the discount; a buyer whose marginal funding costs more should pay on day 30 --- never on day 15,
which gives up the discount without using the credit.

\pitfalls
\begin{itemize}
  \item Computing the rate on the full invoice ($500/50{,}000 = 1\%$): the amount borrowed is the discounted price.
  \item Counting 30 days of credit instead of 20; the first 10 days are free under both choices.
  \item Annualising by simple interest ($1.0101\% \times 18.25 = 18.43\%$), which understates the cost.
\end{itemize}

\refbib Ross Cap.~20, \S20.2 ``Condi\c{c}\~oes de venda'', ``Custo de cr\'edito'', p.~709. No printed answer.

% =====================================================================
\exsec{4. Kyoto Boy: collection period with a discount --- Ross 20.3, p.~733}
%% BEGIN-STATEMENT 4
\stmtbegin
\textbf{Source:} Ross, Cap. 20, Exercise 20.3, p.~733 (PDF 757) · Answer in book: no

\textbf{PMR e contas a receber} {[}OA1{]} A Kyoto Boy Ltda. vende previsões de rentabilidade de
títulos mobiliários japoneses. Suas condições de crédito são de pagamento em 30 dias com opção de
desconto de 2\% para pagamento até o 10º dia. Com base na experiência, 65\% de todos os clientes
aceitarão o desconto.

\begin{enumerate}
\def\labelenumi{\alph{enumi}.}
\item
  Qual é o prazo médio de recebimento da empresa?
\item
  Se a empresa vender 1.300 previsões todos os meses, a \$1.700 cada uma, qual será o valor médio de
  contas a receber no balanço patrimonial?
\end{enumerate}
\stmtend
%% END-STATEMENT

\asked (a) The average collection period (PMR) when 65\% of customers pay on day 10 (taking
the 2\% discount) and the rest pay on day 30. (b) The average receivables balance for 1{,}300
sales a month at \$1{,}700 each.

\approach The PMR is the payment day averaged over customers, weighted by the share that pays on
each day (Ross, ``O desconto financeiro e o PMR'', p.~710). Receivables then follow from
$AR = \text{average daily sales} \times PMR$ (Ross, \S20.1, Eq.\ for ``Contas a receber'').

\notation
\begin{itemize}
  \item $w = 0.65$: share of customers taking the discount; $PMR$: average collection period (days).
  \item $S_m$, $S_a$: monthly and annual sales (\$, before discounts); $AR$: average receivables (\$).
\end{itemize}
\assume Customers pay exactly on day 10 or day 30; sales are measured before discounts (as in
Ross's own example, p.~710); daily sales $=$ annual sales$/365$.

\step{1}{(a) Weighted collection period}
\begin{align}
  PMR &= w \cdot 10 + (1-w)\cdot 30 \label{r5be4-pmr}\\
  &= (0.65)(10) + (0.35)(30) \notag\\
  &= 6.5 + 10.5 \notag\\
  &= 17 \ \text{days}. \notag
\end{align}

\answer{(a)} $PMR = 17$ days.

\step{2}{(b) Annual and daily sales}
\begin{align}
  S_m &= (1{,}300)(1{,}700) = 2{,}210{,}000, \notag\\
  S_a &= 12\,(2{,}210{,}000) = 26{,}520{,}000, \notag\\
  S_a/365 &= (26{,}520{,}000)/365 = 72{,}657.53 \ \text{per day}. \notag
\end{align}

\step{3}{(b) Receivables}
\begin{align}
  AR &= (S_a/365)\times PMR = (72{,}657.53)(17) = 1{,}235{,}178.08. \label{r5be4-ar}
\end{align}

\answer{(b)} Average receivables $\approx \$1{,}235{,}178$.

\intuition The discount pulls two-thirds of the customers forward by 20 days, so the firm carries
only 17 days of sales instead of 30. The price is the 2\% given away on 65\% of sales; whether it
pays depends on the value of the 13 days of financing released, the subject of Exercise 3 seen
from the seller's side.

\pitfalls
\begin{itemize}
  \item Averaging the days without weights ($(10+30)/2 = 20$).
  \item Multiplying the PMR by \emph{monthly} sales; with a 30-day month the result would be
        $2{,}210{,}000 \times 17/30 = 1{,}252{,}333.33$, a different (and inconsistent) day count.
  \item Netting the discount out of sales without saying so; it changes the answer by about 1\%.
\end{itemize}

\refbib Ross Cap.~20, \S20.1 ``Cr\'edito e contas a receber'' and \S20.2, ``O desconto financeiro
e o PMR'', p.~710. No printed answer.

% =====================================================================
\exsec{5. Happy Valley Homecare: inventory days --- Berk 19.17, p.~589}
%% BEGIN-STATEMENT 5
\stmtbegin
\textbf{Source:} Berk, Ch. 19, Exercise 17, p.~589 (PDF 620) · Answer in book: no

Happy Valley Homecare Suppliers, Inc.~(HVHS), had \$20 million in sales in 2010. Its cost of goods
sold was \$8 million, and its average inventory balance was \$2,000,000.

\begin{enumerate}
\def\labelenumi{\alph{enumi}.}
\item
  Calculate the average number of days inventory outstanding ratios for HVHS.
\item
  The average days of inventory in the industry is 73 days. By how much would HVHS reduce its
  investment in inventory if it could improve its inventory days to meet the industry average?
\end{enumerate}
\stmtend
%% END-STATEMENT

\asked (a) How many days of COGS HVHS holds in inventory. (b) How much less inventory it would
carry at the industry's 73 days.

\approach Inventory days $=$ inventory$/$(COGS$/365$) (Berk, \S19.1 and \S19.5). Inverting the
ratio, the inventory a firm must carry at a target number of days is target days $\times$ daily
COGS; the difference with the current balance is cash released once.

\notation
\begin{itemize}
  \item $C = 8{,}000{,}000$: COGS; $\overline{INV} = 2{,}000{,}000$: average inventory (\$).
  \item $D_{INV}$: inventory days; $D^* = 73$: industry days; $INV^*$: inventory at $D^*$ (\$).
\end{itemize}

\step{1}{(a) Daily COGS and inventory days}
\begin{align}
  C/365 &= (8{,}000{,}000)/365 = 21{,}917.81, \notag\\
  D_{INV} &= \overline{INV}/(C/365) = (2{,}000{,}000)/(21{,}917.81) = 91.25 \ \text{days}. \label{r5be5-dio}
\end{align}

\answer{(a)} $D_{INV} = 91.25$ days.

\step{2}{(b) Inventory at the industry average}
\begin{align}
  INV^* &= D^*\,(C/365) = (73)(21{,}917.81) = 1{,}600{,}000, \notag\\
  \overline{INV} - INV^* &= (2{,}000{,}000) - (1{,}600{,}000) = 400{,}000. \label{r5be5-red}
\end{align}
Check: $(91.25 - 73)\times 21{,}917.81 = (18.25)(21{,}917.81) = 400{,}000$.

\answer{(b)} Inventory investment falls by \$400{,}000.

\intuition Each day of inventory ties up one day of COGS. Cutting 18.25 days frees \$400{,}000
once (a one-time cash inflow), plus the carrying cost on it every year after. Sales (\$20 million)
play no role: inventory is at cost.

\pitfalls
\begin{itemize}
  \item Using sales instead of COGS ($2{,}000{,}000/(20{,}000{,}000/365) = 36.5$ days), which
        makes HVHS look far better than the industry.
  \item Reporting the saving as annual; it is a one-time release of capital.
\end{itemize}

\refbib Berk Ch.~19, \S19.1, p.~567; \S19.5 ``Inventory Management'', p.~580. No printed answer.

% =====================================================================
\exsec{6. The Greek Connection: NWC and the CCC --- Berk 19.6, p.~587}
%% BEGIN-STATEMENT 6
\stmtbegin
\textbf{Source:} Berk, Ch. 19, Exercise 6, p.~587 (PDF 618) · Answer in book: no

The Greek Connection had sales of \$32 million in 2010, and a cost of goods sold of \$20 million. A
simplified balance sheet for the firm appears below:

\begin{enumerate}
\def\labelenumi{\alph{enumi}.}
\item
  Calculate The Greek Connection's net working capital in 2010.
\item
  Calculate the cash conversion cycle of The Greek Connection in 2010.
\item
  The industry average accounts receivable days is 30 days. What would the cash conversion cycle for
  The Greek Connection have been in 2010 had it matched the industry average for accounts receivable
  days?
\end{enumerate}

\textbf{The Greek Connection --- Balance Sheet as of December 31, 2010 (thousands of dollars)}

{\def\LTcaptype{none} % do not increment counter
\begin{longtable}[]{@{}llll@{}}
\toprule\noalign{}
Assets & & Liabilities and Equity & \\
\midrule\noalign{}
\endhead
\bottomrule\noalign{}
\endlastfoot
Cash & \$2,000 & Accounts payable & \$1,500 \\
Accounts receivable & 3,950 & Notes payable & 1,000 \\
Inventory & 1,300 & Accruals & 1,220 \\
Total current assets & \$7,250 & Total current liabilities & \$3,720 \\
Net plant, property, and equipment & 8,500 & Long-term debt & \$3,000 \\
& & Total liabilities & \$6,720 \\
Total assets & \$15,750 & Common equity & \$9,030 \\
& & Total liabilities and equity & \$15,750 \\
\end{longtable}
}
\stmtend
%% END-STATEMENT

\asked (a) Net working capital from the balance sheet. (b) The cash conversion cycle. (c) The
CCC if receivables days had been the industry's 30 days.

\approach Berk defines net working capital as current assets minus current liabilities (\S19.1);
the CCC is built as in Exercise 1. Part (c) changes only the receivables component, so the CCC
moves one-for-one with receivables days.

\notation
\begin{itemize}
  \item $CA$, $CL$: current assets and liabilities (\$ thousand); $NWC = CA - CL$.
  \item $S = 32{,}000$, $C = 20{,}000$: sales and COGS (\$ thousand); days ratios as in Exercise 1.
\end{itemize}

\step{1}{(a) Net working capital}
\begin{align}
  CA &= 2{,}000 + 3{,}950 + 1{,}300 = 7{,}250, \notag\\
  CL &= 1{,}500 + 1{,}000 + 1{,}220 = 3{,}720, \notag\\
  NWC &= CA - CL = (7{,}250) - (3{,}720) = 3{,}530. \label{r5be6-nwc}
\end{align}

\answer{(a)} $NWC = \$3{,}530$ thousand (\$3.53 million).

\noindent Remark: the \emph{operating} part, which excludes cash and interest-bearing notes, is
$AR + INV - AP - \text{Accruals} = 3{,}950 + 1{,}300 - 1{,}500 - 1{,}220 = 2{,}530$ thousand.

\step{2}{(b) Daily flows and days ratios}
\begin{align}
  S/365 &= (32{,}000)/365 = 87.671, \qquad C/365 = (20{,}000)/365 = 54.795, \notag\\
  D_{AR} &= (3{,}950)/(87.671) = 45.05 \ \text{days}, \notag\\
  D_{INV} &= (1{,}300)/(54.795) = 23.73 \ \text{days}, \notag\\
  D_{AP} &= (1{,}500)/(54.795) = 27.38 \ \text{days}. \notag
\end{align}

\step{3}{(b) Cash conversion cycle}
\begin{align}
  CCC &= D_{INV} + D_{AR} - D_{AP} = (23.73) + (45.05) - (27.38) = 41.40 \ \text{days}. \label{r5be6-ccc}
\end{align}

\answer{(b)} $CCC = 41.40$ days.

\step{4}{(c) Receivables at 30 days}
Replace $D_{AR}$ by 30 in \eqref{r5be6-ccc}:
\begin{align}
  CCC' &= (23.73) + (30) - (27.38) = 26.35 \ \text{days}, \label{r5be6-cccc}\\
  CCC - CCC' &= (41.40) - (26.35) = 15.05 \ \text{days}. \notag
\end{align}
The receivables balance consistent with 30 days and the cash it would release are
\begin{align}
  AR' &= (30)(87.671) = 2{,}630.14, \notag\\
  AR - AR' &= (3{,}950) - (2{,}630.14) = 1{,}319.86 \ \text{thousand}. \notag
\end{align}

\answer{(c)} $CCC' = 26.35$ days (15.05 days shorter), releasing about \$1.32 million of receivables.

\intuition Greek Connection's customers take 45 days against an industry norm of 30: the firm is
lending to its customers for an extra 15 days. Matching the norm would cut the CCC by a third
and release \$1.32 million --- more than a third of its NWC --- that could repay the notes payable.

\pitfalls
\begin{itemize}
  \item Computing receivables days on COGS ($3{,}950/54.795 = 72$ days).
  \item Leaving notes payable out of current liabilities in (a); NWC in Berk's definition includes them.
  \item In (c), recomputing all three components; only receivables days change.
\end{itemize}

\refbib Berk Ch.~19, \S19.1 ``Overview of Working Capital'', pp.~566--567; \S19.3 ``Receivables
Management'', p.~576. No printed answer.

% =====================================================================
\exsec{7. Terms 3/15, net 40 --- Berk 19.14, p.~588}
%% BEGIN-STATEMENT 7
\stmtbegin
\textbf{Source:} Berk, Ch. 19, Exercise 14, p.~588 (PDF 619) · Answer in book: no

Your firm purchases goods from its supplier on terms of 3/15, net 40.

\begin{enumerate}
\def\labelenumi{\alph{enumi}.}
\item
  What is the effective annual cost to your firm if it chooses not to take the discount and makes
  its payment on day 40?
\item
  What is the effective annual cost to your firm if it chooses not to take the discount and makes
  its payment on day 50?
\end{enumerate}
\stmtend
%% END-STATEMENT

\asked The effective annual cost of skipping a 3\% discount (a) when paying on the due date, day
40, and (b) when stretching payment to day 50.

\approach Skipping the discount borrows the discounted price, $97\%$ of the invoice, from day 15
until the payment day, at a cost of $3\%$ of the invoice. The rate per period is $3/97$; the
period is the number of days of credit actually used, and the EAR compounds that period over a
365-day year (Berk, Example 19.3 and Example 19.7, ``stretching the accounts payable'').

\notation
\begin{itemize}
  \item $d = 0.03$: discount; $t_d = 15$: discount deadline; $t$: actual payment day.
  \item $r = d/(1-d)$: rate per credit period; $n = t - t_d$: days of credit (days).
\end{itemize}

\step{1}{Rate per credit period}
On a \$100 invoice, pay \$97 on day 15 or \$100 later: \$3 of interest on \$97 borrowed,
\begin{align}
  r = \frac{d}{1-d} = \frac{(0.03)}{(0.97)} = 3.0928\%. \label{r5be7-r}
\end{align}

\step{2}{(a) Pay on day 40}
Credit period $n = 40 - 15 = 25$ days; there are $365/25 = 14.6$ such periods per year:
\begin{align}
  \text{EAR}_{40} &= (1+r)^{365/25} - 1 \notag\\
  &= (1.030928)^{14.6} - 1 \notag\\
  &= 56.00\%. \label{r5be7-a}
\end{align}

\answer{(a)} Effective annual cost $= 56.00\%$.

\step{3}{(b) Pay on day 50}
Now $n = 50 - 15 = 35$ days and $365/35 = 10.4286$ periods:
\begin{align}
  \text{EAR}_{50} &= (1.030928)^{10.4286} - 1 \notag\\
  &= 37.39\%. \label{r5be7-b}
\end{align}

\answer{(b)} Effective annual cost $= 37.39\%$.

\intuition The interest charge is fixed (\$3 per \$97), so the longer the firm keeps the money,
the lower the rate per day: stretching from 40 to 50 days cuts the cost from 56\% to 37\%. But
both are far above normal bank rates, so the first-best is usually to borrow from the bank and
take the discount. Stretching has hidden costs the formula does not see --- late-payment penalties,
a damaged credit rating, suppliers that refuse future credit.

\pitfalls
\begin{itemize}
  \item Using 40 or 50 days as the period instead of the days beyond the discount date (25, 35).
  \item Using $3/100$ instead of $3/97$.
  \item Reporting the APR $3.0928\% \times 14.6 = 45.15\%$ as the cost; Berk's ``Common Mistake''
        box (p.~572) warns that the EAR is the right comparison.
\end{itemize}

\refbib Berk Ch.~19, \S19.2 ``Trade Credit'', Cost of Trade Credit, pp.~571--572 (Example 19.3);
\S19.4 ``Payables Management'', Stretching Accounts Payable, p.~579 (Example 19.7). No printed answer.

% =====================================================================
\exsec{8. Toyota vs.\ Apple: operating and cash conversion cycles --- Higgins 2.13, p.~75}
%% BEGIN-STATEMENT 8
\stmtbegin
\textbf{Source:} Higgins, Ch. 2, Exercise 13, p.~75 (PDF 96) · Answer in book: yes, p.~404

The following table presents selected 2020 annual income statement items and balance sheet items for
Toyota Motor and Apple. All figures are in millions. Use the information to answer the questions
that follow.

{\def\LTcaptype{none} % do not increment counter
\begin{longtable}[]{@{}lll@{}}
\toprule\noalign{}
& Toyota & Apple \\
\midrule\noalign{}
\endhead
\bottomrule\noalign{}
\endlastfoot
Sales & ¥ 29,929,992 & \$ 274,515 \\
Cost of goods sold & 24,522,364 & 169,559 \\
Accounts receivable & 8,709,065 & 16,120 \\
Inventory & 2,434,918 & 4,061 \\
Accounts payable & 2,434,180 & 42,296 \\
\end{longtable}
}

\begin{enumerate}
\def\labelenumi{\alph{enumi}.}
\item
  Calculate and interpret the length of the operating cycle for Toyota and Apple.
\item
  Calculate and interpret the length of the cash conversion cycle for Toyota and Apple.
\item
  What are the advantages and disadvantages of Apple's way of running operations compared to
  Toyota's way?
\end{enumerate}
\stmtend
%% END-STATEMENT

\asked (a) The operating cycle of each firm (inventory days $+$ collection period) and what it
means. (b) The cash conversion cycle (operating cycle $-$ payables period) and what it means. (c)
The pros and cons of Apple's way of running operations relative to Toyota's.

\approach Higgins (Ch.~2, ``Operating and Cash Conversion Cycles'') defines days inventory
outstanding and the payables period on COGS, the collection period on sales, all with 365 days.
Because each firm's ratios are in its own currency per day, the yen/dollar difference cancels
inside each ratio and the days are directly comparable.

\notation
\begin{itemize}
  \item $D_{INV}$, $D_{AR}$, $D_{AP}$: days inventory outstanding, collection period, payables period (days).
  \item $OC = D_{INV} + D_{AR}$; $CCC = OC - D_{AP}$ (days). Figures in millions (\yen{} or \$).
\end{itemize}

\step{1}{Daily flows}
\begin{align}
  \text{Toyota: } C/365 &= (24{,}522{,}364)/365 = 67{,}184.56, & S/365 &= (29{,}929{,}992)/365 = 82{,}000.0, \notag\\
  \text{Apple: } C/365 &= (169{,}559)/365 = 464.55, & S/365 &= (274{,}515)/365 = 752.10. \notag
\end{align}

\step{2}{(a) Inventory days and collection period}
\begin{align}
  D_{INV}^{T} &= (2{,}434{,}918)/(67{,}184.56) = 36.24, & D_{AR}^{T} &= (8{,}709{,}065)/(82{,}000.0) = 106.21, \notag\\
  D_{INV}^{A} &= (4{,}061)/(464.55) = 8.74, & D_{AR}^{A} &= (16{,}120)/(752.10) = 21.43. \notag
\end{align}

\step{3}{(a) Operating cycles}
\begin{align}
  OC^{T} &= (36.24) + (106.21) = 142.45 \ \text{days}, \label{r5be8-oct}\\
  OC^{A} &= (8.74) + (21.43) = 30.18 \ \text{days}. \label{r5be8-oca}
\end{align}

\answer{(a)} Toyota $142.45$ days; Apple $30.18$ days. On average Toyota needs about 142 days from
acquiring inventory to collecting the customer's cash, Apple only 30. Toyota's cycle is dominated
by receivables: 106 days, which reflects dealer financing and Toyota's captive finance arm
(customer loans and leases sit in receivables).

\step{4}{(b) Payables periods}
\begin{align}
  D_{AP}^{T} &= (2{,}434{,}180)/(67{,}184.56) = 36.23 \ \text{days}, \notag\\
  D_{AP}^{A} &= (42{,}296)/(464.55) = 91.05 \ \text{days}. \notag
\end{align}

\step{5}{(b) Cash conversion cycles}
\begin{align}
  CCC^{T} &= (142.45) - (36.23) = 106.22 \ \text{days}, \label{r5be8-ccct}\\
  CCC^{A} &= (30.18) - (91.05) = -60.87 \ \text{days}. \label{r5be8-ccca}
\end{align}

\answer{(b)} Toyota $106.22$ days; Apple $-60.87$ days. Toyota pays for its inputs about 106 days
before it collects, so it must finance that gap; Apple collects from customers about 61 days
\emph{before} it pays its suppliers, so its suppliers finance its operations and more. In balance
sheet terms, Apple's trade working capital $AR + INV - AP = 16{,}120 + 4{,}061 - 42{,}296 =
-22{,}115$ million is a source of funds, while Toyota's, $8{,}709{,}065 + 2{,}434{,}918 -
2{,}434{,}180 = 8{,}709{,}803$ million yen, is a large use.

\step{6}{(c) Advantages and disadvantages of Apple's model}
\begin{itemize}
  \item \emph{Advantages.} (i) No need to finance working capital; growth generates cash instead of
        consuming it, because each extra dollar of sales adds more payables than receivables and
        inventory. (ii) Little inventory means little obsolescence risk in a fast-changing product
        line. (iii) Bargaining power over suppliers and outsourced assembly lets Apple push the
        holding of inventory and the waiting for payment upstream.
  \item \emph{Disadvantages and risks.} (i) Thin inventory risks stock-outs and lost sales when
        demand spikes or a supply chain is disrupted. (ii) Demanding quick payment can drive
        customers (carriers, retailers) away. (iii) Paying slowly can strain or alienate
        suppliers, raise the prices they charge, or create fragility if a supplier fails. (iv) The
        model depends on market power; a firm without Apple's brand cannot copy it.
  \item \emph{Why Toyota differs.} Selling cars on credit (dealer floor-plan and consumer
        financing) is part of Toyota's product offer; the long receivables period is a business
        choice funded by its financial-services arm, not a collection failure.
\end{itemize}

\intuition The CCC tells you who finances whom along the supply chain. A negative cycle is a
quasi-free loan from suppliers, available only to firms with market power and fast inventory
turnover.

\pitfalls
\begin{itemize}
  \item Converting yen to dollars before computing days; it is unnecessary because each ratio is
        within one currency.
  \item Basing the payables period on sales instead of COGS.
  \item Calling Apple's negative CCC ``risk-free''; it comes with supplier and stock-out risks.
\end{itemize}

\bookans Higgins p.~404: DIO 36.2 and 8.7 days; collection period 106.2 and 21.4 days; operating
cycle 142.4 and 30.1 days; payables 36.2 and 91.0 days; CCC 106.2 and $-60.9$ days; (c) the same
advantages and risks. Matches the book's answer; the book adds rounded components (e.g.\ $36.2 +
106.2 = 142.4$, $8.7 + 21.4 = 30.1$), which explains the differences in the second decimal.

\refbib Higgins Ch.~2, ``Operating and Cash Conversion Cycles'', p.~47.

% =====================================================================
\exsec{9. Hand-to-Mouth: forgone discount vs.\ two bank loans --- Berk 20.8, p.~613}
%% BEGIN-STATEMENT 9
\stmtbegin
\textbf{Source:} Berk, Ch. 20, Exercise 8, p.~613 (PDF 644) · Answer in book: no

The Hand-to-Mouth Company needs a \$10,000 loan for the next 30 days. It is trying to decide which
of three alternatives to use:

\begin{itemize}
\tightlist
\item
  \textbf{Alternative A:} Forgo the discount on its trade credit agreement that offers terms of
  2/10, net 30.
\item
  \textbf{Alternative B:} Borrow the money from Bank A, which has offered to lend the firm \$10,000
  for 30 days at an APR of 12\%. The bank will require a (no-interest) compensating balance of 5\%
  of the face value of the loan and will charge a \$100 loan origination fee, which means
  Hand-to-Mouth must borrow even more than the \$10,000.
\item
  \textbf{Alternative C:} Borrow the money from Bank B, which has offered to lend the firm \$10,000
  for 30 days at an APR of 15\%. The loan has a 1\% loan origination fee.
\end{itemize}

Which alternative is the cheapest source of financing for Hand-to-Mouth?
\stmtend
%% END-STATEMENT

\asked Which of three ways to get \$10{,}000 of usable funds for about 30 days is cheapest: (A)
skipping a 2/10, net 30 discount; (B) a 12\% APR loan with a 5\% non-interest compensating
balance and a \$100 fee; (C) a 15\% APR loan with a 1\% origination fee.

\approach Compare effective annual rates computed on the cash the firm can actually use. A fee
paid up front and a compensating balance both reduce usable proceeds while interest is charged on
the face value, so the loan must be grossed up until usable funds equal \$10{,}000 (Berk,
\S20.3, ``Common Loan Stipulations and Fees''). The trade-credit cost is as in Exercise 7.

\notation
\begin{itemize}
  \item $L$: face value borrowed (\$); $i$: interest for the 30-day loan (\$); $r$: 30-day rate on usable funds.
  \item EAR $= (1+r)^{365/n} - 1$ with $n$ the length of the credit period (days).
\end{itemize}
\assume A 30-day loan is one month: interest $= (\text{APR}/12)\,L$, paid at maturity. The
compensating balance is applied to the repayment at maturity (Berk's Timmons example, p.~603).
All three rates are annualised with 365 days, $365/n$ periods, so they are comparable.

\step{1}{Alternative A: forgo the discount}
Paying \$98 on day 10 or \$100 on day 30 means borrowing \$98 for $n = 20$ days:
\begin{align}
  r_A &= (2)/(98) = 2.0408\%, \notag\\
  \text{EAR}_A &= (1.020408)^{365/20} - 1 = (1.020408)^{18.25} - 1 = 44.59\%. \label{r5be9-a}
\end{align}

\step{2}{Alternative B: gross up the face value}
Usable proceeds $= L - 0.05L - 100 = 10{,}000$, so
\begin{align}
  0.95\,L &= 10{,}000 + 100 = 10{,}100, \notag\\
  L &= (10{,}100)/(0.95) = 10{,}631.58. \notag
\end{align}
The compensating balance is $0.05\,L = 531.58$. Interest and net repayment at day 30:
\begin{align}
  i_B &= (0.12/12)(10{,}631.58) = 106.32, \notag\\
  \text{Net repayment} &= L + i_B - 0.05L = (10{,}631.58) + (106.32) - (531.58) = 10{,}206.32, \notag\\
  r_B &= (10{,}206.32)/(10{,}000) - 1 = 2.0632\%, \notag\\
  \text{EAR}_B &= (1.020632)^{365/30} - 1 = (1.020632)^{12.1667} - 1 = 28.21\%. \label{r5be9-b}
\end{align}

\step{3}{Alternative C: gross up for the 1\% fee}
Usable proceeds $= 0.99\,L = 10{,}000$:
\begin{align}
  L &= (10{,}000)/(0.99) = 10{,}101.01 \quad (\text{fee } 0.01L = 101.01), \notag\\
  i_C &= (0.15/12)(10{,}101.01) = 126.26, \notag\\
  \text{Repayment} &= (10{,}101.01) + (126.26) = 10{,}227.27, \notag\\
  r_C &= (10{,}227.27)/(10{,}000) - 1 = 2.2727\%, \notag\\
  \text{EAR}_C &= (1.022727)^{12.1667} - 1 = 31.45\%. \label{r5be9-c}
\end{align}

\step{4}{Ranking}
$\text{EAR}_B = 28.21\% < \text{EAR}_C = 31.45\% < \text{EAR}_A = 44.59\%$. With 12 monthly periods
instead of $365/30$ the loan EARs are $27.77\%$ and $30.95\%$; the ranking is unchanged.

\answer{\unskip} Alternative B (Bank A) is cheapest, at an EAR of about 28\%; forgoing the discount is by
far the most expensive (44.6\%).

\intuition The stated APRs (12\% vs.\ 15\%) understate both loans' costs, because the borrower pays
interest on money it cannot use (the fee, the idle balance). Bank A's lower APR survives these
add-ons. Trade credit that ``costs nothing'' in cash outlay is the most expensive source here: 2\%
for only 20 days. Note also that alternative A finances 20 days, not the 30 needed --- another
reason to prefer a bank loan and take the discount.

\pitfalls
\begin{itemize}
  \item Comparing APRs (12\%, 15\%) directly, ignoring the fee and the compensating balance.
  \item Borrowing only \$10{,}000 in B and C: usable proceeds would then fall short of the need.
  \item Forgetting that the compensating balance is returned at maturity (it offsets the repayment).
\end{itemize}

\refbib Berk Ch.~20, \S20.3 ``Short-Term Financing with Bank Loans'', Common Loan Stipulations and
Fees, pp.~602--603; Ch.~19, \S19.2, pp.~571--572. No printed answer.

% =====================================================================
\exsec{10. Pe\c{c}as A\'ereas: one-off vs.\ repeat order --- Ross 20.9, pp.~733--734}
%% BEGIN-STATEMENT 10
\stmtbegin
\textbf{Source:} Ross, Cap. 20, Exercise 20.9, pp.~733--734 (PDF 757) · Answer in book: no

\textbf{Avaliação da política de crédito} {[}OA2{]} A Peças Aéreas é um atacadista que armazena
componentes de motor e equipamento de testes para a indústria de aviões comerciais. Um novo cliente
fez um pedido de oito motores à turbina de alta derivação que aumentam a economia de combustível. O
custo variável é de \$1,6 milhão por unidade, e o preço a prazo é de \$1,87 milhão cada. O crédito é
concedido por um período e, com base na experiência histórica, o pagamento de aproximadamente 1 de
cada 200 pedidos nunca é recebido. O retorno exigido é de 2,9\% por período.

\begin{enumerate}
\def\labelenumi{\alph{enumi}.}
\item
  Assumindo que é um pedido único, ele deve ser atendido? O cliente não comprará se o crédito não
  for concedido.
\item
  Qual é a probabilidade de inadimplência no ponto de equilíbrio da parte (a)?
\item
  Suponha que os clientes que paguem tornem-se clientes e façam o mesmo pedido a cada período, para
  sempre. Assuma também que esses clientes nunca deixem de pagar. O pedido deve ser atendido? Qual é
  a probabilidade de inadimplência no ponto de equilíbrio?
\item
  Descreva em termos gerais o motivo pelo qual as condições de crédito serão mais liberais quando
  pedidos repetidos forem uma possibilidade.
\end{enumerate}

\paragraph{Context needed to solve it}

\emph{From Ross, Cap. 20, ``Quando o crédito deve ser concedido?'' (Uma venda única; Negócios
repetidos), pp.~717--718 (PDF 741--742) --- the one-period credit model and the NPV formulas behind
parts (a)--(d):}

\textbf{Uma venda única} Começaremos avaliando o caso mais simples. Um cliente novo deseja comprar
uma unidade a prazo, a um preço \(P\) por unidade. Se o crédito for negado, o cliente não fará a
compra.

Além disso, assumimos que, se o crédito for concedido, em um mês o cliente fará ou não o pagamento.
A probabilidade de o segundo evento ocorrer é \(\pi\). Nesse caso, a probabilidade (\(\pi\)) pode
ser interpretada como a porcentagem dos novos clientes que não pagarão. Nossa empresa não tem
clientes antigos. Portanto, a rigor, essa é uma venda única. Por fim, o retorno exigido sobre as
contas a receber é \(R\) por mês, e o custo variável é \(v\) por unidade.

A análise aqui é objetiva. Se a empresa negar o crédito, o fluxo de caixa incremental será zero. Se
ela conceder o crédito, gastará \(v\) (o custo variável) neste mês e esperará receber \((1 - \pi)P\)
no mês seguinte. O VPL da concessão de crédito será:

\[\text{VPL} = -v + (1 - \pi)P/(1 + R) \qquad [20.8]\]

{[}\ldots{]} Observe que dividimos por \((1 + R)\) aqui em vez de \(R\) porque, agora, assumimos que
essa é uma transação única.

Nosso exemplo ilustra um ponto importante. Ao conceder crédito a um novo cliente, uma empresa
arrisca seu custo variável (\(v\)). Isso permite obter o preço total (\(P\)). Portanto, para um
cliente novo, o crédito pode ser concedido mesmo que a probabilidade de inadimplência seja alta. Por
exemplo, nesse caso, a probabilidade no ponto de equilíbrio pode ser determinada definindo o VPL
como zero e calculando \(\pi\): {[}\ldots{]}

\textbf{Negócios repetidos} Um segundo e importante fator a ser lembrado é a possibilidade de
negócios repetidos. Podemos ilustrar isso ampliando nosso exemplo da venda única. Fazemos em
pressuposto importante aqui: um cliente novo que paga na primeira vez permanecerá cliente para
sempre e nunca deixará de pagar.

Se a empresa conceder crédito, ela gastará \(v\) neste mês. No mês seguinte, ela nada terá se o
cliente não pagar ou terá \(P\) se o cliente pagar. Se o cliente pagar, ele comprará outra unidade a
prazo, e a empresa gastará \(v\) novamente. Portanto, a entrada de caixa líquida do mês será
\(P - v\). A cada mês subsequente, esse mesmo \(P - v\) ocorrerá à medida que o cliente pagar o
pedido do mês anterior e fizer outro pedido.

Concluímos da nossa discussão que, em um mês, a empresa receberá \$0 com probabilidade \(\pi\). No
entanto, com probabilidade \((1 - \pi)\), a empresa terá um cliente novo permanente. O valor do
cliente novo é igual ao valor presente de \((P - v)\) a cada mês para sempre:

\[VP = (P - v)/R\]

Então, o VPL da concessão de crédito é:

\[\text{VPL} = -v + (1 - \pi)(P - v)/R \qquad [20.9]\]

{[}\ldots{]}
\stmtend
%% END-STATEMENT

\asked (a) Whether to grant one period of credit on a one-off order of eight engines with a 0.5\%
default chance. (b) The default probability that makes that NPV zero. (c) The same two questions
when a customer who pays becomes a permanent, never-defaulting customer. (d) Why repeat business
makes credit terms more liberal.

\approach Ross's one-period credit model (\S20.4, context block in the list): granting credit
risks the variable cost now against the expected price next period. For a one-off sale,
$\text{NPV} = -v + (1-\pi)P/(1+R)$ [20.8]. With repeat business a paying customer is worth a
perpetuity of margins, $\text{NPV} = -v + (1-\pi)(P - v)/R$ [20.9]. Break-even probabilities set
each NPV to zero.

\notation
\begin{itemize}
  \item $v = 1.6$: variable cost per engine; $P = 1.87$: credit price per engine (\$ million).
  \item $\pi = 1/200 = 0.005$: default probability; $R = 0.029$: required return per period.
  \item $\pi^*$: break-even default probability. The order is 8 engines.
\end{itemize}

\step{1}{(a) NPV per engine, one-off sale}
\begin{align}
  \text{NPV} &= -v + \frac{(1-\pi)P}{1+R} \label{r5be10-npv}\\
  &= -1.6 + \frac{(0.995)(1.87)}{1.029} \notag\\
  &= -1.6 + \frac{1.86065}{1.029} \notag\\
  &= -1.6 + 1.808212 \notag\\
  &= 0.208212 \ \text{per engine}. \notag
\end{align}
For the order: $8 \times 0.208212 = 1.665695$ (\$ million).

\answer{(a)} Yes: NPV $\approx \$1.67$ million for the eight-engine order ($\$208{,}212$ per engine).

\step{2}{(b) Break-even default probability}
Set \eqref{r5be10-npv} to zero and solve for $\pi$:
\begin{align}
  0 &= -v + (1-\pi^*)P/(1+R) \notag\\
  (1-\pi^*)P &= v(1+R) \notag\\
  1-\pi^* &= \frac{v(1+R)}{P} = \frac{(1.6)(1.029)}{1.87} = \frac{1.6464}{1.87} = 0.880428 \notag\\
  \pi^* &= 1 - 0.880428 = 0.119572. \label{r5be10-pib}
\end{align}

\answer{(b)} $\pi^* \approx 11.96\%$.

\step{3}{(c) Repeat business}
A paying customer generates $P - v = 1.87 - 1.6 = 0.27$ per engine every period forever:
\begin{align}
  \text{NPV} &= -v + \frac{(1-\pi)(P-v)}{R} \label{r5be10-npvc}\\
  &= -1.6 + \frac{(0.995)(0.27)}{0.029} \notag\\
  &= -1.6 + \frac{0.26865}{0.029} \notag\\
  &= -1.6 + 9.263793 \notag\\
  &= 7.663793 \ \text{per engine}, \notag
\end{align}
i.e.\ $8 \times 7.663793 = 61.310345$ (\$ million) for the order. Break-even:
\begin{align}
  (1-\pi^*)(P - v)/R &= v \notag\\
  1-\pi^* &= \frac{vR}{P - v} = \frac{(1.6)(0.029)}{0.27} = \frac{0.0464}{0.27} = 0.171852 \notag\\
  \pi^* &= 1 - 0.171852 = 0.828148. \label{r5be10-pic}
\end{align}

\answer{(c)} Yes: NPV $\approx \$61.31$ million for the order; break-even default probability $\approx 82.81\%$.

\step{4}{(d) Why repeat business liberalises credit}
\begin{itemize}
  \item In a one-off sale the firm risks $v$ to win one payment $P$; the upside is a single margin.
  \item With repeat business, the first credit sale is a screening device: risking $v$ once buys,
        with probability $1-\pi$, a customer whose value is the perpetuity $(P - v)/R$. That prize
        is many times $P$ (here $9.31$ vs.\ $1.87$ per engine), so the firm tolerates a much higher
        default rate on new customers ($82.8\%$ vs.\ $12.0\%$).
  \item Defaulters reveal themselves after one period and are dropped, so the loss is capped at
        one unit of variable cost, while good customers keep paying. Credit is cheap
        information about who is a good customer.
\end{itemize}

\intuition Granting credit to a new customer is an option-like investment: limited downside ($v$
once), large upside (a stream of margins). The more valuable the relationship, the looser the
credit standard.

\pitfalls
\begin{itemize}
  \item Discounting the perpetuity in (c) by $1+R$ or using $P$ instead of $P - v$: the customer
        pays $P$ but the firm spends $v$ again on each repeat order.
  \item Computing the NPV with $\pi$ applied to the cost as well as the price; the cost is spent
        for sure.
  \item Forgetting to multiply by the eight engines when reporting the order's NPV.
\end{itemize}

\refbib Ross Cap.~20, \S20.4 ``Pol\'itica de cr\'edito \'otima'' / ``Quando o cr\'edito deve ser
concedido?'' (Uma venda \'unica; Neg\'ocios repetidos), pp.~717--718. No printed answer.

% =====================================================================
\exsec{11. Lince S/A: seasonal cash budget and short-term plan --- Ross 18.15, p.~657}
%% BEGIN-STATEMENT 11
\stmtbegin
\textbf{Source:} Ross, Cap. 18, Exercise 18.15, p.~657 (PDF 681) · Answer in book: no

\textbf{Cálculo do orçamento de caixa} {[}OA3{]} A Lince S/A estimou as seguintes vendas (em
milhões) para os próximos quatro trimestres:

{\def\LTcaptype{none} % do not increment counter
\begin{longtable}[]{@{}lllll@{}}
\toprule\noalign{}
& T1 & T2 & T3 & T4 \\
\midrule\noalign{}
\endhead
\bottomrule\noalign{}
\endlastfoot
Vendas & \$210 & \$180 & \$245 & \$280 \\
\end{longtable}
}

As vendas do primeiro trimestre do ano seguinte são projetadas em \$240 milhões. As contas a receber
no início do ano são de \$68 milhões. A empresa tem um prazo médio de recebimento de 45 dias.

A Lince S/A compra dos fornecedores em um trimestre o equivalente a 45\% das vendas previstas para o
próximo trimestre, e os fornecedores normalmente são pagos em 36 dias. Os salários, impostos e
outras despesas custam cerca de 25\% das vendas. Os juros e dividendos são de \$12 milhões por
trimestre.

Para o segundo trimestre, está planejado um desembolso de capital de \$80 milhões. Por último, o
saldo de caixa inicial é de \$64 milhões, e a empresa deseja manter um saldo mínimo de \$30 milhões.

\begin{enumerate}
\def\labelenumi{\alph{enumi}.}
\tightlist
\item
  Complete o orçamento de caixa para a Lince S/A:
\end{enumerate}

\textbf{LINCE S/A --- Orçamento de caixa (em milhões)}

{\def\LTcaptype{none} % do not increment counter
\begin{longtable}[]{@{}lllll@{}}
\toprule\noalign{}
& T1 & T2 & T3 & T4 \\
\midrule\noalign{}
\endhead
\bottomrule\noalign{}
\endlastfoot
Meta para o saldo de caixa & \$30 & & & \\
Entrada de caixa líquida & & & & \\
Saldo de caixa final & & & & \\
Saldo de caixa mínimo & 30 & & & \\
Superávit (déficit) cumulativo & & & & \\
\end{longtable}
}

\begin{enumerate}
\def\labelenumi{\alph{enumi}.}
\setcounter{enumi}{1}
\tightlist
\item
  Assuma que a Lince S/A possa financiar a curto prazo todos os fundos necessários à taxa de 3\% por
  trimestre e possa investir todo o excedente de fundos em instrumentos financeiros de curto prazo a
  uma taxa de 2\% por trimestre. Prepare um plano financeiro de curto prazo preenchendo a
  programação a seguir. Qual é o custo de caixa líquido (juros totais pagos menos os lucros totais
  de investimento) para o ano?
\end{enumerate}

\textbf{LINCE S/A --- Plano financeiro de curto prazo (em milhões)}

{\def\LTcaptype{none} % do not increment counter
\begin{longtable}[]{@{}lllll@{}}
\toprule\noalign{}
& T1 & T2 & T3 & T4 \\
\midrule\noalign{}
\endhead
\bottomrule\noalign{}
\endlastfoot
Meta para o saldo de caixa & \$64 & & & \\
Entrada de caixa líquida & & & & \\
Novos investimentos de curto prazo & & & & \\
Lucros dos investimentos de curto prazo & & & & \\
Investimentos de curto prazo resgatados & & & & \\
Novo empréstimo de curto prazo & & & & \\
Juros sobre o empréstimo de curto prazo & & & & \\
Amortização do empréstimo de curto prazo & & & & \\
Saldo de caixa final & & & & \\
Saldo de caixa mínimo & & & & \\
Sobra (falta) acumulada & 30 & & & \\
Investimentos de curto prazo inicial & & & & \\
Investimentos de curto prazo final & & & & \\
Dívida de curto prazo inicial & & & & \\
Dívida de curto prazo final & & & & \\
\end{longtable}
}
\stmtend
%% END-STATEMENT

\asked (a) Fill in the quarterly cash budget: net cash inflow, ending cash and the cumulative
surplus over the \$30 million minimum. (b) Build the short-term financial plan in which surpluses
earn 2\% per quarter and deficits are borrowed at 3\% per quarter, and give the net cash cost
(interest paid minus investment income) for the year.

\approach Ross's cash budget (\S18.4): collections follow from the collection period, payments
from the payables period, both as fractions of a 90-day quarter; net cash inflow $=$ collections
$-$ disbursements. The short-term plan (\S18.6) then invests any cash above the minimum and
borrows to cover any shortfall, with interest computed on the beginning-of-quarter balances.

\notation
\begin{itemize}
  \item $S_q$: sales in quarter $q$; $Pur_q = 0.45\,S_{q+1}$: purchases in $q$ (\$ million).
  \item $Col_q$, $Pay_q$: collections and payments to suppliers; $NCF_q$: net cash inflow.
  \item $I_q$: short-term investments at the end of $q$; $Inc_q = 0.02\,I_{q-1}$: investment income.
\end{itemize}
\assume A quarter has 90 days. Collection period 45 days: $45/90 = 1/2$ of a quarter's sales is
collected in the quarter, the other half in the next; the \$68 million opening receivables are
collected in Q1. Suppliers paid in 36 days: $36/90 = 0.4$ of a quarter's purchases is paid the
next quarter, $0.6$ in the same quarter. Last year's Q4 purchases were $0.45 \times 210 = 94.5$
(45\% of this year's Q1 sales). The book's first row, \pt{``Meta para o saldo de caixa \$30''}, is
read as the beginning cash balance, which the statement fixes at \$64 million (the \$30 is the
minimum balance, and the part (b) table itself starts at \$64).

\step{1}{Collections}
\begin{align}
  Col_1 &= 68 + \tfrac12(210) = 68 + 105 = 173, \notag\\
  Col_2 &= \tfrac12(210) + \tfrac12(180) = 105 + 90 = 195, \notag\\
  Col_3 &= \tfrac12(180) + \tfrac12(245) = 90 + 122.5 = 212.5, \notag\\
  Col_4 &= \tfrac12(245) + \tfrac12(280) = 122.5 + 140 = 262.5. \notag
\end{align}

\step{2}{Purchases and payments to suppliers}
Purchases: $Pur_1 = 0.45(180) = 81$, $Pur_2 = 0.45(245) = 110.25$, $Pur_3 = 0.45(280) = 126$,
$Pur_4 = 0.45(240) = 108$. Payments $Pay_q = 0.4\,Pur_{q-1} + 0.6\,Pur_q$:
\begin{align}
  Pay_1 &= 0.4(94.5) + 0.6(81) = 37.8 + 48.6 = 86.4, \notag\\
  Pay_2 &= 0.4(81) + 0.6(110.25) = 32.4 + 66.15 = 98.55, \notag\\
  Pay_3 &= 0.4(110.25) + 0.6(126) = 44.1 + 75.6 = 119.7, \notag\\
  Pay_4 &= 0.4(126) + 0.6(108) = 50.4 + 64.8 = 115.2. \notag
\end{align}

\step{3}{Disbursements and net cash inflow}
Wages, taxes and other: $0.25 S_q = 52.5,\ 45,\ 61.25,\ 70$; interest and dividends 12 each
quarter; capital spending 80 in Q2.
\begin{center}
\begin{tabular}{lrrrr}
\toprule
 & Q1 & Q2 & Q3 & Q4\\
\midrule
Payments to suppliers & 86.40 & 98.55 & 119.70 & 115.20\\
Wages, taxes, other & 52.50 & 45.00 & 61.25 & 70.00\\
Interest and dividends & 12.00 & 12.00 & 12.00 & 12.00\\
Capital spending & 0.00 & 80.00 & 0.00 & 0.00\\
\midrule
Total disbursements & 150.90 & 235.55 & 192.95 & 197.20\\
Collections & 173.00 & 195.00 & 212.50 & 262.50\\
\midrule
Net cash inflow & 22.10 & $-40.55$ & 19.55 & 65.30\\
\bottomrule
\end{tabular}
\end{center}
For example $NCF_2 = 195 - (98.55 + 45 + 12 + 80) = 195 - 235.55 = -40.55$.

\step{4}{(a) Cash budget}
Ending cash $=$ beginning cash $+$ $NCF_q$; cumulative surplus $=$ ending cash $- 30$:
\begin{center}
\begin{tabular}{lrrrr}
\toprule
 & Q1 & Q2 & Q3 & Q4\\
\midrule
Beginning cash balance & 64.00 & 86.10 & 45.55 & 65.10\\
Net cash inflow & 22.10 & $-40.55$ & 19.55 & 65.30\\
Ending cash balance & 86.10 & 45.55 & 65.10 & 130.40\\
Minimum cash balance & 30.00 & 30.00 & 30.00 & 30.00\\
Cumulative surplus (deficit) & 56.10 & 15.55 & 35.10 & 100.40\\
\bottomrule
\end{tabular}
\end{center}

\answer{(a)} Ending cash 86.10, 45.55, 65.10, 130.40; cumulative surplus 56.10, 15.55, 35.10,
100.40 (\$ million). There is never a deficit.

\step{5}{(b) Short-term plan}
Hold exactly \$30 million; invest the excess; earn $2\%$ on the investments held at the start of
each quarter. Q1: invest $64 + 22.10 - 30 = 56.10$. Q2: income $0.02(56.10) = 1.122$; the deficit
$-40.55 + 1.122 = -39.428$ is covered by selling investments, so $I_2 = 56.10 - 39.428 = 16.672$.
Q3: income $0.02(16.672) = 0.33344$; invest $19.55 + 0.33344 = 19.88344$; $I_3 = 36.55544$. Q4:
income $0.02(36.55544) = 0.73111$; invest $65.30 + 0.73111 = 66.03111$; $I_4 = 102.58655$.
\begin{center}\small
\begin{tabular}{lrrrr}
\toprule
 & Q1 & Q2 & Q3 & Q4\\
\midrule
Beginning cash balance & 64.000 & 30.000 & 30.000 & 30.000\\
Net cash inflow & 22.100 & $-40.550$ & 19.550 & 65.300\\
New short-term investments & 56.100 & 0.000 & 19.883 & 66.031\\
Income from short-term investments & 0.000 & 1.122 & 0.333 & 0.731\\
Short-term investments sold & 0.000 & 39.428 & 0.000 & 0.000\\
New short-term borrowing & 0 & 0 & 0 & 0\\
Interest on short-term borrowing & 0 & 0 & 0 & 0\\
Short-term borrowing repaid & 0 & 0 & 0 & 0\\
Ending cash balance & 30.000 & 30.000 & 30.000 & 30.000\\
Minimum cash balance & 30.000 & 30.000 & 30.000 & 30.000\\
Cumulative surplus (deficit) & 0.000 & 0.000 & 0.000 & 0.000\\
Beginning short-term investments & 0.000 & 56.100 & 16.672 & 36.555\\
Ending short-term investments & 56.100 & 16.672 & 36.555 & 102.587\\
Beginning short-term debt & 0 & 0 & 0 & 0\\
Ending short-term debt & 0 & 0 & 0 & 0\\
\bottomrule
\end{tabular}
\end{center}
Net cash cost $=$ interest paid $-$ investment income:
\begin{align}
  0 - (0 + 1.122 + 0.33344 + 0.73111) = -2.18655. \label{r5be11-cost}
\end{align}

\answer{(b)} No borrowing is ever needed; the net cash cost is $-\$2.19$ million, i.e.\ a net
investment \emph{income} of about \$2.19 million for the year.

\intuition Lince's seasonal pattern (big Q2 capital spending, strong Q4 collections) is fully
covered by the opening cash cushion of \$64 million: the Q2 dip is met by selling part of what was
invested in Q1. The cash budget is the tool that reveals this timing; annual totals alone would
hide the Q2 squeeze. With a lower opening balance the same plan would show Q2 borrowing at 3\%.

\pitfalls
\begin{itemize}
  \item Basing purchases on the current quarter's sales instead of next quarter's, or forgetting
        last year's Q4 purchases (94.5) in the Q1 payment.
  \item Collecting 45/90 of the \emph{previous} quarter's sales only, i.e.\ swapping the two halves.
  \item Computing investment income on end-of-quarter balances (it is earned on what was invested
        during the quarter, the beginning balance).
\end{itemize}

\refbib Ross Cap.~18, \S18.4 ``O or\c{c}amento de caixa'', pp.~639--642; \S18.6 ``Um plano
financeiro de curto prazo'', pp.~649--650. No printed answer.

% =====================================================================
\exsec{12. BCD Manufacturing: tax shield vs.\ lost supplier credit --- Grinblatt 17.5, p.~619}
%% BEGIN-STATEMENT 12
\stmtbegin
\textbf{Source:} Grinblatt \& Titman, Ch. 17, Exercise 17.5, p.~619 (PDF 632) · Answer in book: no

BCD Manufacturing is considering repurchasing 40 percent of its common stock. Management estimates
the tax savings from such a move to be \$48 million, based on the addition of \$1 billion of debt at
a rate of 12 percent with a 40 percent marginal tax rate. However, the company's suppliers are
unhappy with the decision and are threatening to revoke the company's net-30 day credit terms, which
will cost the firm an additional 2 percent on its \$1.5 billion inventory. Should management go
ahead with the repurchase? Why or why not?
\stmtend
%% END-STATEMENT

\asked Whether BCD should borrow \$1 billion to buy back 40\% of its stock, when the debt's tax
saving is \$48 million a year but suppliers threaten to withdraw net-30 terms, which would cost an
extra 2\% a year on \$1.5 billion of inventory; and why.

\approach Grinblatt \& Titman's stakeholder theory of capital structure (\S17.1): leverage creates
tax benefits but imposes costs through non-financial stakeholders --- here suppliers who fear
being unpaid and react by tightening credit. Compare the two flows on the same (after-tax, annual)
basis, then ask what else the suppliers' reaction signals.

\notation
\begin{itemize}
  \item $D = 1{,}000$: new debt; $r_D = 0.12$: interest rate; $\tau = 0.40$: marginal tax rate (\$ million).
  \item $INV = 1{,}500$: inventory; $c = 0.02$: extra annual financing cost if trade credit is lost.
\end{itemize}
\assume The \$48 million and the 2\% cost are annual and recur as long as the debt is in place;
the 2\% is a pre-tax, tax-deductible financing cost (it replaces trade credit with costlier funds).

\step{1}{Annual tax saving}
\begin{align}
  \text{Interest} &= r_D D = (0.12)(1{,}000) = 120, \notag\\
  \text{Tax saving} &= \tau\, r_D D = (0.40)(120) = 48. \label{r5be12-ts}
\end{align}

\step{2}{Annual cost of losing trade credit}
\begin{align}
  \text{Pre-tax cost} &= c\, INV = (0.02)(1{,}500) = 30, \notag\\
  \text{After-tax cost} &= (1-\tau)(30) = (0.60)(30) = 18. \label{r5be12-sc}
\end{align}

\step{3}{Net annual effect}
\begin{align}
  48 - 18 = 30 \ \text{per year} \quad (\text{even with no tax deduction: } 48 - 30 = 18 > 0). \notag
\end{align}
In present value, with perpetual debt, the tax shield is $\tau D = (0.40)(1{,}000) = 400$, while
the after-tax supplier cost discounted at 12\% is $18/0.12 = 150$; the net is $400 - 150 = 250$.

\step{4}{Qualitative judgement}
The narrow arithmetic favours the repurchase, but the suppliers' reaction is a warning sign:
\begin{itemize}
  \item Suppliers are a stakeholder group that bears costs in distress; if they react, customers
        (who worry about warranties and spare parts) and employees (who worry about their jobs) may
        too, and those indirect distress costs are not in the \$18 million.
  \item The \$48 million is not riskless: it is earned only if BCD has taxable income to shelter;
        higher leverage raises the chance of losses and of financial distress.
  \item Personal taxes on interest reduce the net tax advantage of debt (Miller), so \$48 million
        overstates the gain to shareholders.
  \item The supplier credit could perhaps be preserved by negotiation (collateral, shorter terms,
        a smaller repurchase).
\end{itemize}

\answer{\unskip} On the numbers given, go ahead: the tax saving (\$48 million) exceeds the after-tax cost
of lost trade credit (\$18 million) by \$30 million a year (\$250 million in PV for perpetual
debt). But management should first weigh the indirect costs the supplier threat foreshadows; if
other stakeholders react similarly, the decision can flip.

\intuition Capital structure is not only a deal between shareholders and lenders. Firms whose
suppliers, customers or employees make firm-specific investments pay for leverage through worse
terms from them. The optimal debt ratio is lower when those stakeholders are important and
easily spooked.

\pitfalls
\begin{itemize}
  \item Comparing the \$48 million tax saving with the \$120 million of interest: interest replaces
        dividends to equity; only the tax saving is a gain to the firm.
  \item Comparing an after-tax gain with a pre-tax cost without saying so.
  \item Stopping at the arithmetic and ignoring what the suppliers' reaction implies about distress costs.
\end{itemize}

\refbib Grinblatt \& Titman Ch.~17, \S17.1 ``The Stakeholder Theory of Capital Structure'',
pp.~597--600. No printed answer.

% =====================================================================
\exsec{13. International Motor Corporation: CCC and payables --- Berk 19.15, p.~588}
%% BEGIN-STATEMENT 13
\stmtbegin
\textbf{Source:} Berk, Ch. 19, Exercise *15, p.~588 (PDF 619) · Answer in book: no

Use the financial statements supplied below and on the next page for International Motor Corporation
(IMC) to answer the following questions:

\begin{enumerate}
\def\labelenumi{\alph{enumi}.}
\item
  Calculate the cash conversion cycle for IMC for both 2009 and 2010. What change has occurred, if
  any? All else being equal, how does this change affect IMC's need for cash?
\item
  IMC's suppliers offer terms of net 30. Does it appear that IMC is doing a good job of managing its
  accounts payable?
\end{enumerate}

\textbf{International Motor Corporation --- Income Statement (in millions) for the Years Ending
December 31}

{\def\LTcaptype{none} % do not increment counter
\begin{longtable}[]{@{}lll@{}}
\toprule\noalign{}
& 2009 & 2010 \\
\midrule\noalign{}
\endhead
\bottomrule\noalign{}
\endlastfoot
Sales & \$60,000 & \$75,000 \\
Cost of goods sold & 52,000 & 61,000 \\
Gross profit & \$8,000 & \$14,000 \\
Selling, general, and administrative expenses & 6,000 & 8,000 \\
Operating profit & \$2,000 & \$6,000 \\
Interest expense & 1,400 & 1,300 \\
Earnings before taxes & \$600 & \$4,700 \\
Taxes & 300 & 2,350 \\
Earnings after taxes & \$300 & \$2,350 \\
\end{longtable}
}

\textbf{International Motor Corporation --- Balance Sheet (in millions) as of December 31}

{\def\LTcaptype{none} % do not increment counter
\begin{longtable}[]{@{}llllll@{}}
\toprule\noalign{}
Assets & 2009 & 2010 & Liabilities & 2009 & 2010 \\
\midrule\noalign{}
\endhead
\bottomrule\noalign{}
\endlastfoot
Cash & \$3,080 & \$6,100 & Accounts payable & \$3,600 & \$4,600 \\
Accounts Receivable & 2,800 & 6,900 & Notes payable & 1,180 & 1,250 \\
Inventory & 6,200 & 6,600 & Accruals & 5,600 & 6,211 \\
Total current assets & \$12,080 & \$19,600 & Total current liabilities & \$10,380 & \$12,061 \\
Net plant, property, and equipment & \$23,087 & \$20,098 & Long-term debt & \$6,500 & \$7,000 \\
& & & Total liabilities & \$16,880 & \$19,061 \\
Total assets & \$35,167 & \$39,698 & \textbf{Equity} & & \\
& & & Common stock & \$2,735 & \$2,735 \\
& & & Retained earnings & \$15,552 & \$17,902 \\
& & & Total equity & \$18,287 & \$20,637 \\
& & & Total liabilities and equity & \$35,167 & \$39,698 \\
\end{longtable}
}
\stmtend
%% END-STATEMENT

\asked (a) IMC's cash conversion cycle in 2009 and 2010, what changed, and what that implies for
its cash needs. (b) Whether, with suppliers' terms of net 30, IMC manages its payables well.

\approach The CCC as in Exercise 1, computed for each year from that year's flows and year-end
balances, then decomposed to see which component moved. For (b), compare payables days with the
30 days allowed: with net terms and no discount, paying early has no benefit (Berk, \S19.4,
``Determining Accounts Payable Days Outstanding'').

\notation
\begin{itemize}
  \item $D_{INV}^{y}$, $D_{AR}^{y}$, $D_{AP}^{y}$, $CCC^{y}$ for year $y \in \{2009, 2010\}$ (days).
  \item Flows and balances in \$ million.
\end{itemize}

\step{1}{Daily flows}
\begin{align}
  2009:&\quad C/365 = (52{,}000)/365 = 142.47, \quad S/365 = (60{,}000)/365 = 164.38, \notag\\
  2010:&\quad C/365 = (61{,}000)/365 = 167.12, \quad S/365 = (75{,}000)/365 = 205.48. \notag
\end{align}

\step{2}{(a) 2009}
\begin{align}
  D_{INV} &= (6{,}200)/(142.47) = 43.52, \quad D_{AR} = (2{,}800)/(164.38) = 17.03, \notag\\
  D_{AP} &= (3{,}600)/(142.47) = 25.27, \notag\\
  CCC^{2009} &= (43.52) + (17.03) - (25.27) = 35.28 \ \text{days}. \label{r5be13-09}
\end{align}

\step{3}{(a) 2010}
\begin{align}
  D_{INV} &= (6{,}600)/(167.12) = 39.49, \quad D_{AR} = (6{,}900)/(205.48) = 33.58, \notag\\
  D_{AP} &= (4{,}600)/(167.12) = 27.52, \notag\\
  CCC^{2010} &= (39.49) + (33.58) - (27.52) = 45.55 \ \text{days}. \label{r5be13-10}
\end{align}

\step{4}{(a) What changed}
$CCC^{2010} - CCC^{2009} = 45.55 - 35.28 = 10.26$ days. By component: inventory $-4.03$ days,
receivables $+16.55$ days, payables $+2.26$ days (which shortens the CCC). The deterioration is
entirely receivables: had 2010 receivables days stayed at 17.03, receivables would have been
$17.03 \times 205.48 = 3{,}500$ instead of 6{,}900, i.e.\ \$3{,}400 million less. The extra
10.26 days of the CCC, valued at 2010 daily sales, tie up about $10.26 \times 205.48 \approx
\$2{,}109$ million.

\answer{(a)} CCC $= 35.28$ days (2009) and $45.55$ days (2010), a rise of $10.26$ days driven by
slower collections (receivables days $17.03 \to 33.58$). All else equal IMC needs more cash (on
the order of \$2 billion) to finance working capital.

\step{5}{(b) Payables against net 30}
Payables days were $25.27$ (2009) and $27.52$ (2010), both below the 30 days suppliers allow. With
net-30 terms there is no discount for paying early, so every day paid before day 30 is free
financing given up. Paying on day 30 in 2010 would mean payables of $30 \times 167.12 =
5{,}013.70$ instead of 4{,}600, i.e.\ \$413.70 million more interest-free credit
($30 \times 142.47 - 3{,}600 = 673.97$ in 2009).

\answer{(b)} No: IMC pays in about 25--28 days although it may take 30 at no cost; it could raise
roughly \$414 million of free financing (2010) by paying on the due date. It has improved slightly
from 2009, and it should not stretch beyond 30 days.

\intuition IMC's working capital got worse because customers pay slower --- possibly a deliberate
loosening of credit to push the 25\% sales growth, possibly a collection problem; the aging
schedule would tell. On the payables side it leaves free money on the table. Both point to
working capital as a cheap source of the cash its growth needs.

\pitfalls
\begin{itemize}
  \item Using one year's flows with the other year's balances.
  \item Advising IMC to stretch payables beyond 30 days; the costless target is exactly the due date.
  \item Reading the higher CCC as ``more cash on hand''; it means more cash is \emph{needed}.
\end{itemize}

\refbib Berk Ch.~19, \S19.1, p.~567; \S19.4 ``Payables Management'', pp.~578--579. No printed answer.

% =====================================================================
\exsec{14. Senna S/A: switch-to-credit NPV and break-even price --- Ross 20.18 (with 20.14), pp.~734--735}
%% BEGIN-STATEMENT 14
\stmtbegin
\textbf{Source:} Ross, Cap. 20, Exercise 20.18 (DESAFIO), p.~735 (PDF 759); Part 0 = Exercise 20.14,
p.~734 (PDF 758) · Answer in book: no

\begin{enumerate}
\def\labelenumi{\arabic{enumi}.}
\setcounter{enumi}{17}
\tightlist
\item
  \textbf{\emph{Markup} de crédito} {[}OA2{]} No Problema 14, qual é o preço por unidade no ponto de
  equilíbrio que deve ser cobrado de acordo com a nova política de crédito? Assuma que, na nova
  política, o número de vendas seja de 4.100 unidades e que todos os outros valores permaneçam
  iguais.
\end{enumerate}

\paragraph{Context needed to solve it}

\textbf{Part 0} --- \emph{Source: Ross, Cap. 20, Exercise 20.14, p.~734 (PDF 758) · Answer in book:
no}

\begin{enumerate}
\def\labelenumi{\arabic{enumi}.}
\setcounter{enumi}{13}
\tightlist
\item
  \textbf{Avaliação da política de crédito} {[}OA2{]} A Senna S/A está pensando em mudar sua
  política atual de vendas somente à vista. A nova condição seria para vendas com pagamento em um
  período pelo valor integral da compra. Com base nas seguintes informações, determine se a empresa
  deve ou não continuar com a proposta. O retorno exigido é de 2,5\% por período.
\end{enumerate}

{\def\LTcaptype{none} % do not increment counter
\begin{longtable}[]{@{}lll@{}}
\toprule\noalign{}
& Política atual & Nova política \\
\midrule\noalign{}
\endhead
\bottomrule\noalign{}
\endlastfoot
Preço por unidade & \$91 & \$94 \\
Custo por unidade & \$47 & \$47 \\
Vendas unitárias por mês & 3.850 & 3.940 \\
\end{longtable}
}
\stmtend
%% END-STATEMENT

\asked Part 0 (Problem 14): whether to switch from cash-only sales at \$91 to one-period credit
at \$94 with sales rising from 3{,}850 to 3{,}940 units. Problem 18: the price per unit under the
new policy that makes the switch's NPV zero when new-policy sales are 4{,}100 units.

\approach Ross's NPV of switching credit policy (\S20.3): the switch costs this period's cash
receipts under the old policy plus the variable cost of the extra units, $PQ + v(Q' - Q)$ [20.5];
from next period on, the firm gains the difference in monthly margins forever. Ross's formula
[20.6] assumes the same price; with a new price $P'$ the perpetual increment is
$(P'-v)Q' - (P - v)Q$. The break-even price sets the NPV to zero.

\notation
\begin{itemize}
  \item $P = 91$, $P'$: old and new price; $v = 47$: unit cost (\$); $Q = 3{,}850$, $Q'$: old and new units per period.
  \item $R = 0.025$: required return per period.
\end{itemize}

\step{1}{Part 0 --- cost of switching}
\begin{align}
  \text{Cost} &= PQ + v(Q' - Q) \label{r5be14-cost}\\
  &= (91)(3{,}850) + (47)(3{,}940 - 3{,}850) \notag\\
  &= 350{,}350 + 4{,}230 = 354{,}580. \notag
\end{align}

\step{2}{Part 0 --- perpetual benefit and NPV}
\begin{align}
  \Delta CF &= (P' - v)Q' - (P - v)Q \notag\\
  &= (94 - 47)(3{,}940) - (91 - 47)(3{,}850) \notag\\
  &= 185{,}180 - 169{,}400 = 15{,}780 \ \text{per period}, \notag\\
  \text{NPV} &= -354{,}580 + 15{,}780/0.025 = -354{,}580 + 631{,}200 = 276{,}620. \label{r5be14-npv0}
\end{align}
Part 0 answer: switch (NPV $= \$276{,}620$).

\step{3}{Problem 18 --- cost of switching at 4{,}100 units}
\begin{align}
  \text{Cost} &= (91)(3{,}850) + (47)(4{,}100 - 3{,}850) = 350{,}350 + 11{,}750 = 362{,}100. \notag
\end{align}

\step{4}{Problem 18 --- solve NPV $= 0$ for $P'$}
\begin{align}
  0 &= -362{,}100 + \frac{(P' - 47)(4{,}100) - 169{,}400}{0.025} \notag\\
  (P' - 47)(4{,}100) - 169{,}400 &= (0.025)(362{,}100) = 9{,}052.5 \notag\\
  (P' - 47)(4{,}100) &= 169{,}400 + 9{,}052.5 = 178{,}452.5 \notag\\
  P' - 47 &= 178{,}452.5/4{,}100 = 43.525 \notag\\
  P' &= 47 + 43.525 = 90.525. \label{r5be14-p}
\end{align}

\answer{\unskip} Break-even credit price $P' \approx \$90.53$ per unit (and in Part 0 the switch has NPV
$\$276{,}620$).

\intuition At 4{,}100 units the extra volume alone pays for one month's delayed receipts, so the
firm could even charge slightly \emph{less} on credit (\$90.53) than for cash (\$91) and break even.
The credit ``markup'' needed is negative: the volume effect does the work.

\pitfalls
\begin{itemize}
  \item Using the new price in the switching cost: what is lost this period is the old cash
        revenue $PQ$, not $P'Q'$.
  \item Applying Ross's [20.6] with $(P - v)(Q' - Q)$ when the price changes; that formula holds
        only for $P' = P$.
  \item Forgetting the extra variable cost $v(Q' - Q)$ in the switching cost.
\end{itemize}

\refbib Ross Cap.~20, \S20.3 ``An\'alise da pol\'itica de cr\'edito'', pp.~712--714 (Eqs.\
20.2--20.7). No printed answer.

% =====================================================================
\exsec{15. Toys-4-Kids: seasonal vs.\ level production --- Higgins 3.14 (with 3.13), pp.~107--109}
%% BEGIN-STATEMENT 15
\stmtbegin
\textbf{Source:} Higgins, Ch. 3, Exercise 14, pp.~108--109 (PDF 129--130); Part 0 = Exercise 13,
pp.~107--108 (PDF 128--129) · Answer in book: no for 14; yes for 13 (Part 0), p.~408

\begin{enumerate}
\def\labelenumi{\arabic{enumi}.}
\setcounter{enumi}{13}
\tightlist
\item
  Continuing with Toys-4-Kids, introduced in the preceding problem, the company's production manager
  has argued for years that it is inefficient to produce on a seasonal basis. She believes the
  company should switch to level production throughout the year, building up finished goods
  inventory in the first two quarters to meet the peak selling needs in the last two. She believes
  the company can reduce its cost of goods sold from 70 percent to 65 percent with level production.
  (Keep in mind that production managers typically want to restrict production to left shoes only so
  as to reduce costs.)
\end{enumerate}

\begin{enumerate}
\def\labelenumi{\alph{enumi}.}
\tightlist
\item
  Prepare a revised pro forma forecast assuming level production. In your forecast, assume that
  quarterly accounts payable under level production equal 10 percent of \emph{average} quarterly
  sales for the year. To estimate quarterly inventory, use the following two formulas:
\end{enumerate}

\[\text{Inventory}_{eoq} = \text{Inventory}_{boq} + \text{Quarterly production} - \text{Quarterly cost of sales}\]

\[\text{Quarterly production} = \text{Annual cost of sales}/4\]

where eoq and boq refer to end of quarter and beginning of quarter, respectively. Please ignore the
effect of increased external financing required on interest expense.

\begin{enumerate}
\def\labelenumi{\alph{enumi}.}
\setcounter{enumi}{1}
\item
  What is the effect of the switch from seasonal to level production on annual profits?
\item
  What effect does the switch have on the company's quarterly ending inventory? On the company's
  quarterly need for external financing?
\item
  Do you think the company will be able to borrow the amount of money required by level production?
  What obsolescence risks does the company incur by building up inventory in anticipation of future
  sales? Might this be a concern to lenders?
\end{enumerate}

\paragraph{Context needed to solve it}

\textbf{Part 0} --- \emph{Source: Higgins, Ch. 3, Exercise 13, pp.~107--108 (PDF 128--129) · Answer
in book: yes, p.~408}

\begin{enumerate}
\def\labelenumi{\arabic{enumi}.}
\setcounter{enumi}{12}
\tightlist
\item
  Toys-4-Kids manufactures plastic toys. Sales and production are highly seasonal. The following is
  a quarterly pro forma forecast indicating external financing needs for 2022. Assumptions are in
  parentheses.
\end{enumerate}

\textbf{Toys-4-Kids --- 2022 Quarterly Pro Forma Forecast (\$ thousands)}

{\def\LTcaptype{none} % do not increment counter
\begin{longtable}[]{@{}lllll@{}}
\toprule\noalign{}
& Qtr 1 & Qtr 2 & Qtr 3 & Qtr 4 \\
\midrule\noalign{}
\endhead
\bottomrule\noalign{}
\endlastfoot
Net sales & 300 & 375 & 3,200 & 5,000 \\
Cost of sales (70 percent of sales) & 210 & 263 & 2,240 & 3,500 \\
Gross profit & 90 & 113 & 960 & 1,500 \\
Operating expenses & 560 & 560 & 560 & 560 \\
Profit before tax & (470) & (448) & 400 & 940 \\
Income taxes & (188) & (179) & 160 & 376 \\
Profit after tax & (282) & (269) & 240 & 564 \\
Cash (minimum balance = \$200,000) & 1,235 & 927 & 200 & 200 \\
Accounts receivable (75 percent of quarterly sales) & 225 & 281 & 2,400 & 3,750 \\
Inventory (12/31/21 balance = \$500,000) & 500 & 500 & 500 & 500 \\
Current assets & 1,960 & 1,708 & 3,100 & 4,450 \\
Net plant \& equipment & 1,000 & 1,000 & 1,000 & 1,000 \\
Total assets & 2,960 & 2,708 & 4,100 & 5,450 \\
Accounts payable (10 percent of quarterly sales) & 30 & 38 & 320 & 500 \\
Accrued taxes (payments quarterly in arrears) & (188) & (179) & 160 & 376 \\
Current liabilities & (158) & (142) & 480 & 876 \\
Long-term debt & 400 & 400 & 400 & 400 \\
Equity (12/31/21 balance = \$3,000,000) & 2,718 & 2,450 & 2,690 & 3,254 \\
Total liabilities and equity & 2,960 & 2,708 & 3,570 & 4,530 \\
External financing required & 0 & 0 & 530 & 920 \\
\end{longtable}
}

\begin{enumerate}
\def\labelenumi{\alph{enumi}.}
\item
  How do you interpret the negative numbers for income taxes in the first two quarters?
\item
  Why are cash balances in the first two quarters greater than the minimum required \$200,000? How
  were these numbers determined?
\item
  How was ``external financing required'' (appearing at the bottom of the forecast) determined?
\item
  Do you think Toys-4-Kids will be able to borrow the external financing required as indicated by
  the forecast?
\end{enumerate}
\stmtend
%% END-STATEMENT

\asked (a) Redo the quarterly pro forma with level production (COGS 65\% of sales, production
$=$ annual COGS$/4$, payables $=$ 10\% of average quarterly sales). (b) The effect on annual
profit. (c) The effect on quarter-end inventory and on external financing. (d) Whether the firm
can borrow that much, the obsolescence risk of building inventory ahead of sales, and whether it
worries lenders.

\approach Percent-of-sales pro forma with external financing required (EFR) as the plug (Higgins
Ch.~3, Problem 13): with cash at its \$200 minimum, EFR $=$ total assets $-$ (liabilities $+$
equity); if EFR is negative, set it to zero and let cash absorb the excess. The forecasts hold
only on each quarter-end date, which is why seasonality matters (Higgins, ``Seasonality'',
pp.~86--87).

\notation
\begin{itemize}
  \item $S_q$: net sales; $C_q = 0.65 S_q$: cost of sales; $Prod = \sum_q C_q/4$: quarterly
        production; $INV_q$: end-of-quarter inventory (\$ thousand).
  \item $\tau = 0.40$: tax rate (from the book's table, $188/470$); $EFR_q$: external financing required.
\end{itemize}
\assume As in Problem 13: operating expenses 560 per quarter, receivables 75\% of quarterly sales,
accrued taxes equal to the quarter's tax (paid in arrears, negative $=$ refund due), net plant
1{,}000, long-term debt 400, opening equity 3{,}000, opening inventory 500, minimum cash 200,
no dividends; interest on new financing ignored, as the problem instructs. Seasonal figures are
recomputed unrounded (the book's table rounds to thousands).

\step{1}{Base case (Problem 13) logic}
In Q1, with cash at 200, assets would be $200 + 225 + 500 + 1{,}000 = 1{,}925$ against liabilities
and equity of $2{,}960$, so EFR would be $-1{,}035$; it is set to 0 and cash becomes $2{,}960 -
1{,}725 = 1{,}235$ (Q2 likewise, EFR $-727$). In Q3 and Q4 EFR is positive: $530.5$ and $920.5$
(the book prints 530 and 920). Annual profit after tax: $(1-0.40)(0.30 \times 8{,}875 - 4 \times
560) = 253.5$.

\step{2}{(a) Level production: cost of sales and production}
Annual sales $\sum S_q = 300 + 375 + 3{,}200 + 5{,}000 = 8{,}875$.
\begin{align}
  C_q &= 0.65 S_q = 195,\ 243.75,\ 2{,}080,\ 3{,}250; \quad \textstyle\sum_q C_q = 5{,}768.75, \notag\\
  Prod &= (5{,}768.75)/4 = 1{,}442.1875 \ \text{per quarter}. \notag
\end{align}

\step{3}{(a) Inventory}
$INV_q = INV_{q-1} + Prod - C_q$, starting from 500:
\begin{align}
  INV_1 &= 500 + 1{,}442.1875 - 195 = 1{,}747.1875, \notag\\
  INV_2 &= 1{,}747.1875 + 1{,}442.1875 - 243.75 = 2{,}945.625, \notag\\
  INV_3 &= 2{,}945.625 + 1{,}442.1875 - 2{,}080 = 2{,}307.8125, \notag\\
  INV_4 &= 2{,}307.8125 + 1{,}442.1875 - 3{,}250 = 500. \notag
\end{align}

\step{4}{(a) Income statement}
Gross profit $S_q - C_q = 105,\ 131.25,\ 1{,}120,\ 1{,}750$; minus 560 gives profit before tax
$-455,\ -428.75,\ 560,\ 1{,}190$; tax at 40\%: $-182,\ -171.5,\ 224,\ 476$; profit after tax
$-273,\ -257.25,\ 336,\ 714$. Equity: $3{,}000 - 273 = 2{,}727$; $2{,}727 - 257.25 = 2{,}469.75$;
$+336 = 2{,}805.75$; $+714 = 3{,}519.75$.

\step{5}{(a) Balance sheet and EFR}
Payables $= 0.10 \times (8{,}875/4) = 0.10 \times 2{,}218.75 = 221.875$ every quarter. With cash at
200, for Q1: assets $200 + 225 + 1{,}747.1875 + 1{,}000 = 3{,}172.1875$; liabilities and equity
$221.875 - 182 + 400 + 2{,}727 = 3{,}166.875$; $EFR_1 = 3{,}172.1875 - 3{,}166.875 = 5.3125$.
The other quarters follow in the same way:
\begin{center}\small
\begin{tabular}{lrrrr}
\toprule
Level production (\$ thousand) & Qtr 1 & Qtr 2 & Qtr 3 & Qtr 4\\
\midrule
Net sales & 300.0 & 375.0 & 3,200.0 & 5,000.0\\
Cost of sales (65\%) & 195.0 & 243.8 & 2,080.0 & 3,250.0\\
Gross profit & 105.0 & 131.3 & 1,120.0 & 1,750.0\\
Operating expenses & 560.0 & 560.0 & 560.0 & 560.0\\
Profit before tax & $-455.0$ & $-428.8$ & 560.0 & 1,190.0\\
Income taxes (40\%) & $-182.0$ & $-171.5$ & 224.0 & 476.0\\
Profit after tax & $-273.0$ & $-257.3$ & 336.0 & 714.0\\
\midrule
Cash & 200.0 & 200.0 & 200.0 & 200.0\\
Accounts receivable (75\%) & 225.0 & 281.3 & 2,400.0 & 3,750.0\\
Inventory & 1,747.2 & 2,945.6 & 2,307.8 & 500.0\\
Net plant \& equipment & 1,000.0 & 1,000.0 & 1,000.0 & 1,000.0\\
Total assets & 3,172.2 & 4,426.9 & 5,907.8 & 5,450.0\\
\midrule
Accounts payable & 221.9 & 221.9 & 221.9 & 221.9\\
Accrued taxes & $-182.0$ & $-171.5$ & 224.0 & 476.0\\
Current liabilities & 39.9 & 50.4 & 445.9 & 697.9\\
Long-term debt & 400.0 & 400.0 & 400.0 & 400.0\\
Equity & 2,727.0 & 2,469.8 & 2,805.8 & 3,519.8\\
Total liabilities and equity & 3,166.9 & 2,920.1 & 3,651.6 & 4,617.6\\
\midrule
External financing required & 5.3 & 1,506.8 & 2,256.2 & 832.4\\
\bottomrule
\end{tabular}
\end{center}

\answer{(a)} EFR under level production: about \$5 thousand, \$1{,}507 thousand, \$2{,}256 thousand
and \$832 thousand in Q1--Q4 (table above).

\step{6}{(b) Annual profit}
Only the cost-of-sales ratio changes, so
\begin{align}
  \Delta PBT &= (0.70 - 0.65)(8{,}875) = 443.75, \notag\\
  \Delta PAT &= (1 - 0.40)(443.75) = 266.25 \quad (519.75 \text{ vs.\ } 253.5). \notag
\end{align}

\answer{(b)} Profit after tax rises by \$266.25 thousand, from \$253.5 to \$519.75 thousand (more
than double).

\step{7}{(c) Inventory and financing}
Inventory rises above the seasonal 500 by $1{,}247.2$, $2{,}445.6$ and $1{,}807.8$ in Q1--Q3 and is
back to 500 in Q4. EFR changes from (0, 0, 530.5, 920.5) to (5.3, 1{,}506.8, 2{,}256.2, 832.4),
i.e.\ by $+5.3$, $+1{,}506.8$, $+1{,}725.7$ and $-88.1$. In Q4 less is needed, because accumulated
extra profit ($+266.25$ of equity) and higher accrued taxes ($+100$) outweigh lower payables
($221.875$ vs.\ 500). The average extra quarter-end financing is $(5.3125 + 1{,}506.75 +
1{,}725.6875 - 88.125)/4 = 3{,}149.625/4 = 787.4$; the pre-tax saving of 443.75 would be wiped out
only at an interest rate of $443.75/787.4 = 56.36\%$ a year --- a rough bound, since it uses
quarter-end balances only.

\answer{(c)} Ending inventory balloons in Q1--Q3 (peak \$2.95 million in Q2) and external
financing is needed from Q1, peaking at about \$2.26 million in Q3 instead of \$0.92 million in Q4;
by year-end the need is slightly lower (\$832 vs.\ \$920 thousand).

\step{8}{(d) Can it borrow, and the obsolescence risk}
\begin{itemize}
  \item \emph{Capacity.} The peak loan (\$2.26 million in Q3) is about $2{,}256/4{,}708 = 47.9\%$
        of receivables plus inventory ($2{,}400 + 2{,}307.8 = 4{,}707.8$); in Q2 the loan is 46.7\% of
        $281.3 + 2{,}945.6 = 3{,}226.9$, but there the collateral is almost all inventory. A
        seasonal, self-liquidating bank line is plausible, but on tougher terms than the base case,
        where loans were under a quarter of receivables (Higgins's answer to 13(d)).
  \item \emph{Obsolescence.} Toys are fashion goods: a line built in Q1--Q2 for Christmas may be
        out of style, mis-forecast or superseded by the time it sells. Unsold goods must be
        discounted or written off, and the planned 5-point margin gain disappears. Level production
        converts a demand forecast into a physical bet made months in advance. (Higgins's aside
        about left shoes is a reminder that what is cheapest to produce is not always what sells.)
  \item \emph{Lenders.} Yes. Inventory is weak collateral: specific to the firm, valued at cost,
        and worth little in liquidation if it is last season's toys. A lender will apply a heavy
        haircut, ask for covenants and monitor sales; repayment depends on the Q3--Q4 selling
        season actually happening. Receivables in Q3--Q4 are better collateral than inventory in
        Q1--Q2.
\end{itemize}

\answer{(d)} Probably yes, as a seasonal line secured by receivables and inventory, but the lender
bears the obsolescence risk of a Q2 balance sheet that is mostly toys not yet sold, and will price
and limit the loan accordingly.

\intuition Level production trades a lower production cost for a larger and earlier investment in
inventory. The profit gain is real and large relative to any plausible interest cost; the hidden
cost is risk --- forecast and obsolescence risk --- which falls on the firm and its lenders.

\pitfalls
\begin{itemize}
  \item Keeping payables at 10\% of each quarter's sales; under level production purchases are
        smooth, hence payables on average sales.
  \item Computing inventory as a percentage of sales instead of with the boq/eoq recursion.
  \item Reading the Q4 EFR (lower than seasonal) as ``level production needs less financing''; the
        peak need is in Q2--Q3.
\end{itemize}

\bookans Problem 14 has no printed answer. For Part 0 (Problem 13), Higgins p.~408: negative taxes
are a reduction of tax liability (or a rebate); cash exceeds the minimum because EFR would be
negative and cash becomes the balancing item; EFR is the balancing item when positive; the firm
should easily borrow, as amounts are less than a quarter of receivables. Step 1 matches this.

\refbib Higgins Ch.~3, ``Pro Forma Statements'' and ``Seasonality'', pp.~86--87; Problems 13--14,
pp.~107--109; answers p.~408.

\end{document}
